\chapter[German]{Deutsch} \label{chap:de}
{\color{intro}\textit{\bt{Die Sprache des deutschen Volkes, schön wie die Klinge eines Dolches. Befürchten Sie keine langen Wörter oder vielfältigen Regeln, aber die Bequemlichkeit, deren Reiz sich nicht leicht ausweichen lässt.}}}

German is a language that has often been insulted by morons that do not understand its usage or beauty, calling it aggressive, making fun of its many features or linking it solely to its history, when in reality, it is not only probably the easiest and most logical one in usage, but also has a history that is not as grim as portrayed when compared to its neighbours, whose misdeeds have just been suppressed because they were on the winning side. With a subtle mix of roughness and smoothness, it is a good starting point for someone trying to learn a new perspective. It also economically opens up basically half of Europe, so there are only advantages to learning it. Alphabet pronunciations are slightly different compared to English, and there exist three characters with double dots (\bt{Umlaut}) unique to German. The convenient part: Save rare exceptions of foreign words, all words are pronounced exactly as they are written. No silent letters or confusing pronunciations. There are only three relevant pronunciation rules that I will mention immediately:
\begin{enumerate}
    \item A \textgerman{\bt{\enquote*{g}}} at the end of a word is pronounced like  \enquote*{sh}
    \item \textgerman{\bt{\enquote*{gt}}} at the end of a word is pronounced like \enquote*{sht}
    \item \textgerman{\bt{\enquote*{ch}}} is pronounced via the throat as  \enquote*{kh}. Sometimes it is also \enquote*{sh}
\end{enumerate}
That's it. Off to a great start, especially considering you already know English with its through, though, thorough, tough\ldots
\section[Noun]{Noun (Nennwort)} \label{sec:denoun}
There is a tri-gender system. Each noun can either be male, female or neutral, and has  singular and plural forms that are modified by  four cases. The problem is that there are multiple patterns for the pluralization of a word for a particular gender. I will list these patterns directly in \autoref{tab:denounT1}, cross reference them with the various words you find as you learn. All  nouns  are in case $1$ - \textgerman{\bt{Nominativ}}, which is also their noun root.
\begin{table}[tbh]
    \centering
    \begin{tabular}{|c|c|c|}
       \hline \bt{Männlich} & \bt{Weiblich} & \bt{Sächlich} \\\hline
        Verlag / Verlage & Schachtel / Schachteln & Gehalt / Gehälter \\\hline
        Körper / Körper & Kraft / Kräfte & Essen / Essen \\\hline
        Sturz / Stürze & Mutter / Mütter & Gericht / Gerichte \\\hline
        Rand / Ränder & Bahn / Bahnen & Gebirge / Gebirge \\\hline
        Boden / Böden & Anwältin / Anwältinnen & \textcolor{red}{Auto / Autos} \\\hline
        Gedanke / Gedanken &  &  \\\hline
        \textcolor{red}{Modus / Modi} & & \\\hline
    \end{tabular}
    \caption{Gender based noun Singular / Plural}
    \label{tab:denounT1}
\end{table}

Any German noun must fit into one of the patterns listed in \autoref{tab:denounT1}. For example,  \bt{Gott / Götter} is covered under \bt{Rand / Ränder} as an \textgerman{\bt{\enquote*{Umlaut + er}}} type plural. The red  entries are foreign (loan) words: their list is not comprehensive. Often, they break traditional German rules. Some abstract nouns cannot be pluralized: \bt{Geduld, Schmach}. Some are only plural: \bt{Semesterferien, Leute}. Some nouns also possess multiple plurals or multiple genders, wherein each gender carries a different meaning:
\begin{example}
    \begin{tabular}{c|c|c}
        \bt{Männlich} & \bt{Gehalt} & Content of something in a mixture \\\hline
        \bt{Sächlich} & \bt{Gehalt} & Salary \\\hline
        \bt{Männlich} & \bt{Kiefer} & Jaw \\\hline
        \bt{Weiblich} & \bt{Kiefer} & Pine tree
    \end{tabular}
\end{example}
Sometimes, the multiple genders do not matter and are  related to dialect differences:
\begin{example}
    \begin{tabular}{c|c|c}
    \bt{Wirrwarr} & \bt{Männlich/Sächlich} & Confusion\\\hline
    \bt{Joghurt} & \bt{Männlich/Weiblich/Sächlich} & Curd
    \end{tabular}
\end{example}
In addition, each noun is attributed an article that represents its gender. Imagine a/an/the but with different forms per case and grammatical gender.

There is a special rule with some male nouns called the N-declination, which causes the forms for all cases (except  $1$) in both singular and plural to become the same as the $1$ plural, which always take the end \bt{-en/n}. The single exception to this rule is the word \bt{Herz} which is the only neutral noun to follow it, that too partially. For neutral nouns in particular, there is an interesting class of words that starts with \bt{ge}. These  words  have some special associated meanings: collective marker or the result of an action:
\begin{example}
    \begin{tabular}{c|c|c}
        \bt{Berg} & \bt{Gebirge} & Mountain vs Mountain range \\\hline
        \bt{Wasser} & \bt{Gewässer} & Water vs Water body\\\hline
        \bt{Malen} & \bt{Gemälde} & To paint vs The result of painting (A painting)\\\hline
        \bt{Stein} & \bt{Gestein} & Stone vs Many stones as a group
    \end{tabular}
\end{example}
These nouns normally do not change form, save for due to the \bt{Dativ, Genetiv} case rules. Note that one cannot conclude in general that a noun starting with \bt{ge}  necessarily falls in this category or is neutral in gender, but such exceptions are rare. Ex: \bt{Geschichte (Weib), Gedanke (Männ), Geisel (Männ/Weib)}. 

An interesting thing  about German words is that they are formed simply, and encode the process of formation itself as a meaning, giving them versatility the likes of which has to be seen to be believed. For example: \bt{Anschlag} is formed of \bt{an} meaning `on/to' and \bt{Schlag} meaning `hit'. So naturally, \bt{Anschlag} is the process of hitting (on/towards) something. This can mean, if performed by one or a group of people, an attack or assassination. If the interpretation is something that is hit on something, then a poster or notice. If considered as the effect of the hit on something, then the sound created by  impact. If perceived as the intention to hit towards  something, then the process of getting into shooting position with a shot based weapon. All of these are valid meanings for this word, and most of German works like this. The consequence is that many times, it is possible to understand new words or make new ones that others may understand, even though a machine translator might get stumped.

All nouns, irrespective of their position in a sentence, must be capitalized: \bt{eine Sache, das Gewässer}
\subsection[Case]{Case (Fall)}
There are four cases in German, not including the trivial $8^{th}$, which means that of the seven we considered, the meanings of three have been compressed into one of the others. Also, each gender has a slightly different pattern of form variation. The case equivalence  along with its corresponding rules are listed in \autoref{tab:denounT2}.
\begin{table}[bth]
    \centering
    \begin{tabular}{|c|c|c|}
        \hline \textit{Official name} & \textit{Case codes} & \textit{Noun forms}\\
        \hline\bt{Nominativ} & $1$ & $\vm{R}$ \\\hline
        \bt{Akkusativ} & $2$ & $\vm{R}$ \\\hline
        \bt{Dativ} & $3,4,5,7$ & $\vm{R}(1)$ / $\vm{R}(1)$ + \bt{e}; $\vm{R}(2)$ + \bt{n}\\\hline
        \bt{Genetiv} & $6$ & $\vm{R}(1)$ / $\vm{R}(1)$ + \bt{e(s)}; $\vm{R}(2)$ \\\hline
    \end{tabular}
    \caption{German case equivalence}
    \label{tab:denounT2}
\end{table}

Male and neutral nouns follow the same pattern, and may have an additional \bt{-e} in their \bt{Dativ} singular. For the \bt{Genetiv} singular, they have either \bt{-es/s} added at the end. The only exception to these rules is when N-declination applies: then all word forms except for $\vm{C}(1,1)$ are equivalent to $\vm{R}(2)$. Female nouns do not follow these rules and just use $\vm{R}(1)$ for $\vm{C}(:,1)$. All nouns must have an additional \bt{n} at the end when building the \bt{Dativ} plural. The only exceptions to these rules are loan words that do not really belong, don't worry about them.  \autoref{tab:denounT3} shows  three nouns: one from each gender, with the male displaying N-declination. 
\begin{table}[bth]
    \centering
    \begin{tabular}{|c|c|c|c|c|c|c|}
        \hline $\Diamond$ & \multicolumn{2}{|c|}{Der Rabe} & \multicolumn{2}{|c|}{Die Bö} & \multicolumn{2}{|c|}{Das Land} \\\hline
        \bt{Nom} & Rabe & Raben & Bö & Böen & Land & Länder \\\hline
        \bt{Akk} & Raben & Raben & Bö & Böen & Land & Länder \\\hline
        \bt{Dat} & Raben & Raben & Bö & Böen & Land/Lande & Ländern\\\hline
        \bt{Gen} & Raben & Raben & Bö & Böen & Landes & Länder \\\hline
    \end{tabular}
    \caption{Examples of some case-multiplicity matrices}
    \label{tab:denounT3}
\end{table}
The easiest test to check whether a word follows the N-declination is to check $\vm{C}(4,1)$. If there is an \bt{-en} type ending, that is a guaranteed N-word. The articles \bt{der/die/das} mentioned in \autoref{tab:denounT3} are typically used to mark genders and themselves have declinations.

A tiny disclaimer: Since female nouns have no word alterations in case $6$, often the preposition \bt{von} is used as an indicator with the meaning `of', especially when no article is used with them. Sometimes, this is also done for other nouns when used without articles. In such formations, the noun follows the \bt{Dativ} case as demanded by \bt{von}.
\begin{example}
    Die Rechte von Frauen, Das Haus von Gott, Der König von Bayern...
\end{example}

\subsection[Article]{Article (Geschlechtswort)}
Correspond to the English (a/an)/the,  and are used to mark the gender of a noun. There exist definite, indefinite and negative articles, whose $\vm{G}_{M}(1,1)$ forms are \bt{der, ein, kein} respectively. All of them decline according to  the noun they are attributed to, even if they stand alone in a sentence as their placeholders. All three types follow similar patterns, whose roots are  $\vm{G}(1,:)$. However, it is most convenient to treat the definite article table \ref{tab:denounT4} as a base for deriving  the others. 
\begin{table}[bth]
    \centering
    \begin{tabular}{|c|c|c|c|c|}
        \hline $\Diamond$ & \bt{Männlich} & \bt{Weiblich} & \bt{Sächlich} & \bt{Mehrzahl} \\\hline
        \bt{Nom} $(1)$ & der & die & das & die \\\hline
        \bt{Akk} $(2)$ & den & die & das & die \\\hline
        \bt{Dat} $(3,4,5,7)$ & dem & der & dem & den \\\hline
        \bt{Gen} $(6)$ & des & der & des & der \\\hline
    \end{tabular}
    \caption{Definite article form variation: $\vm{G}$}
    \label{tab:denounT4}
\end{table}

Some examples of article forms: \bt{keine, keinem, einem, einer, eine} and so on. \bt{Kein} and \bt{ein} decline exactly the same, and have an exceptional configuration for $\vm{G}_{M}(1,1)$ and $\vm{G}_{N}(1 \leftrightarrow2, 1)$, where the forms are \bt{ein/kein} rather than \bt{keiner} for example. A structural difference between them is that \bt{ein} is an analogue to a/an and does not possess a plural form. Since this pattern of variation will repeat often, you should have \autoref{tab:denounT4} memorized. Articles can also be used as  $3^{rd}$ person pronouns in certain contexts.

\subsection[Adjective]{Adjective (Eigenschaftswort)}
Adjectives  also double as adverbs. If a non-capitalized word root  ends in any of the following suffixes,  it is an adjective:
\begin{enumerate}
    \item \bt{-ig}: Highly versatile suffix that can be applied on almost any word \begin{example}
        lauig, traurig, dreckig, heutig, eckig\ldots
    \end{example}
    \item \bt{-bar}: Analogue of the `-able' suffix in English and denotes ability. Applied exclusively on verb roots (VR)\begin{example}
        fahrbar, auffindbar, machbar, verwendbar, kletterbar\ldots
    \end{example}
    \item \bt{-lich}: Normally applied on nouns, although  VRs are also allowed. Often have spelling variations\begin{example}
        klanglich, gefährlich, einheitlich, spärlich, sächlich\ldots
    \end{example}
    \item \bt{-isch}: Implies belonging, or the meaning  `in the context of'. These  usually have many germanized words\begin{example}
        akribisch, städtisch, mechanisch, deutsch, theoretisch\ldots
    \end{example}
    \item Pure variants without any particular pattern\begin{example}
        rot, grausam, lau, öde, frisch\ldots
    \end{example}
    \item All \bt{MdV, MdG} forms
\end{enumerate}
This list is not comprehensive but sufficient.

If embedded in a noun unit, adjectives change form according to noun gender, case and multiplicity. The nature of suffixes used mimics \autoref{tab:denounT4}, so even though \autoref{tab:denounT5} may seem daunting, it contains a lot of redundancy.
\begin{table}[bth]
    \centering
    \begin{tabular}{|c|c|c|c|c|}
        \hline \bt{Fall} & \bt{Männlich} & \bt{Weiblich} & \bt{Sächlich} & \bt{Mehrzahl} \\\hline\hline
        \multicolumn{5}{|c|}{\textit{Definite Article} / Bestimmtes Geschlechtswort}\\\hline
        \bt{Nom} & -e & -e & -e & -en \\\hline
        \bt{Akk} & -en & -e & -e & -en \\\hline
        \bt{Dat} & -en & -en & -en & -en \\\hline
        \bt{Gen} & -en & -en & -en & -en \\\hline\hline
        \multicolumn{5}{|c|}{\textit{Possessive pronoun, Indefinite/Negative article}}\\\hline
        \bt{Nom} & -er & -e & -es & -en \\\hline
        \bt{Akk} & -en & -e & -es & -en \\\hline
        \bt{Dat} & -en & -en & -en & -en \\\hline
        \bt{Gen} & -en & -en & -en & -en \\\hline\hline
        \multicolumn{5}{|c|}{\textit{Without article} / Ohne Geschlechtswort}\\\hline
        \bt{Nom} & -er & -e & -es & -e\\\hline
        \bt{Akk} & -en & -e & -es & -e \\\hline
        \bt{Dat} & -em & -er & -em & -en\\\hline
        \bt{Gen} & -en & -er & -en & -er\\\hline
    \end{tabular}
    \caption{Adjective form variation: $\vm{J}$}
    \label{tab:denounT5}
\end{table}

The idea is that whenever an adjective precedes a noun, the nature of use will determine which table is active. For example: \bt{Der mutig(e) Mann} requires $\vm{J}_{M}(1,1)$ in the definite article table on the adjective \bt{mutig}. The same example using the indefinite article: \bt{Ein mutig(er) Mann}. If an adjective stands alone outside a noun unit, none of this is applicable, and only its root  is used. Such sentences always need the existential verb \bt{sein} to work.
\begin{example}
    \begin{tabular}{c|c}
        Dieses Eis ist lecker & Das leckere Eis...\\\hline
        Das Wetter ist kalt & Das kalte Wetter...
    \end{tabular}
\end{example}
An exception occurs when the adjective (root form: comparatives do not follow this) ends in either \bt{l/r}. In such cases, the last and second last letters are reversed before declining: \bt{dunkel/dunkle, übel/üble}. Do note that I am assuming you have enough of a brain to realize that $\vm{J}_{F}(1,1)$  will not end up as \bt{dunklee} but just \bt{dunkle} because an \bt{e} already exists there. As I said earlier, spelling rules are  for ease of pronunciation, so take the tables with a pinch of salt. If it sounds unnatural, there is probably a minor change somewhere.

In some rare cases, some adjectives are  never declined. Ex: \bt{rosa, lila}. This pattern is supposedly found only in colours ending with vowels (although \bt{blau} declines) and words with foreign origin. To be fair, these are so rare, that you can forget this fact and still never have difficulties. Unless of course, colours are a major part of your life.

The general  comparative and superlative root forms in German are formed similar to English and have the same purpose:
\begin{itemize}[leftmargin=25pt]
    \item Comparative: `root' + \bt{-er} (follows the \bt{l/r} principle during formation)
    \item Superlative:  `root' + \bt{-st}
\end{itemize}
So something like \bt{mutig} (courageous) would end up with comparative \bt{mutiger} and superlative \bt{mutigst}. Another example: \bt{dunkel/dunkler/dunkelst}. These higher order adjectives (\bt{Steigerungsformen}) also follow \autoref{tab:denounT5} if kept in the noun unit. To simplify matters, I will mark the associated cases in brackets for the following subsection.

The comparative is used to compare a property of two nouns specifically, or to compare a noun property to its observed average:
\begin{example}
    \begin{tabular}{m{7cm}|m{7cm}}
 \bt{Er ist mutiger($1$) als er} & He is more courageous than him \\\hline
  \bt{Wir wohnen in einer größeren($7$) Stadt} & We live in a \ut{bigger} city
    \end{tabular}
\end{example}
The second sentence means that the city in which we live is bigger than many others (lies on the larger side), but not necessarily bigger than a particular named example.

The superlative expresses extremes:
\begin{example}
    \begin{tabular}{m{7.5cm}|m{7cm}}
  \bt{Er ist am mutigsten($7$)} & He is the most courageous\\\hline
  \bt{Der mutigste($1$) Ritter ist im Gefecht gestorben} & The bravest knight has died in the battle
    \end{tabular}
\end{example}
In the first example, \bt{am} = \bt{an + dem}, and is a short form. These are explained at the very end in \autoref{sec:despco}, and will be explicitly expanded if used.  Another superlative use case would be phrases like `one of the most popular\ldots' In such cases, you use the indefinite case $1$ article for the noun  being graded, and place $\vm{G}(4,2)$ with the superlative directly after it.  Note that the indefinite articles used here are \bt{einer/eine/eines} for the male/female/neutral cases.
\begin{example}
    \begin{tabular}{m{7.5cm}|m{7.5cm}}
        \bt{Einer der stärksten Männer der Welt ist Yujiro Hanma} & One of the strongest men of the world is Yujiro Hanma \\\hline
        \bt{Eine der größten Städte in Deutschland ist München} & One of the biggest cities in Germany is Munich\\\hline
        \bt{Eines der gerechtesten Räder des Weltalls ist das Schicksalsrad} & One of the most just wheels of the universe is the wheel of fate
    \end{tabular}
\end{example}
Nothing would be natural about natural language without exceptional configurations, which appear for single-syllabic adjectives: a syllable roughly being the extent of a word you can pronounce without significantly changing the state (jaw and tongue position) of your mouth. More precisely, it is a distinct speech unit  around a vowel. In such cases, the comparative receives an \bt{Umlaut} on whichever alphabet can receive it: \bt{ä,ö,ü}.
\begin{example}
    \begin{tabular}{c}
    \bt{groß/größer/größt, rot/röter/rötest}
    \end{tabular}
\end{example}
Then are the adjectives beyond rhyme or reason, but English is no better here:
\begin{example}
    \begin{tabular}{c}
    \bt{gut/besser/best, hoch/höher/höchst, gern/lieber/liebst}
    \end{tabular}
\end{example}
However, not all adjectives have comparative and superlative forms, by virtue of their meaning: \bt{leer} (empty), \bt{schwanger} (pregnant), \bt{tot} (dead). These adjectives are the only ones that do NOT double as adverbs.

\subsection[Pronoun]{Pronoun (Fürwort)}
Pronouns  do not need compulsory capitalization within sentences.  We start with the class of personal pronouns. Since German only has two stages of formality: formal/informal, the formal forms have not been explicitly considered in \autoref{tab:denounT6}.
\begin{table}[bth]
    \centering
    \begin{subtable}{\textwidth}
        \centering
    \begin{tabular}{|c|c|c|c|}
        \hline \multicolumn{2}{|c|}{$1_p$} & \multicolumn{2}{|c|}{$2_p$}\\\hline
        \bt{Ein} &\bt{Mehr} & \bt{Ein} & \bt{Mehr}\\\hline
        ich & wir & du & ihr\\\hline
        mich & uns & dich & euch\\\hline
        mir & uns & dir & euch\\\hline
        mein/meine & unser/unsere & dein/deine & euer/eure\\\hline
    \end{tabular}
    \caption{$K = 1_{p}, 2_{p}$}
\end{subtable}
\begin{subtable}{\textwidth}
    \centering
\begin{tabular}{|c|c|c|c|}
        \hline \multicolumn{4}{|c|}{$3_p$}\\\hline
        \bt{Männ} &\bt{Weib} & \bt{Säch} & \bt{Mehr}\\\hline
        er & sie & es & sie \\\hline
        ihn & sie & es & sie\\\hline
        ihm & ihr & ihm & ihnen\\\hline
        sein/seine & ihr/ihre & sein/seine & ihr/ihre\\\hline
    \end{tabular}
    \caption{$K = 3_{p}$}
\end{subtable}
\caption{Pronoun form variation: $\vm{P}[a,In,K]$}
\label{tab:denounT6}
\end{table}

For  the polite second person, singular and plural for any particular case are the same. Use the plural forms of the given third person, capitalize the first letter, and there you have the polite second person. Important is that in the polite variant, the pronoun MUST be capitalized irrespective of its place in the sentence. This matrix will be stacked in the third dimension above $\vm{P}[a,In,2_p](:,:,0)$ as $\vm{P}[a,In,2_p](:,:,\mathbb{H})$. Case $6$ personal pronouns take on the gender of the noun  being possessed:
\begin{example}
    \begin{tabular}{c|c}
        \bt{Diese Luft ist ihre} & This air is hers\\\hline
        \bt{Die Fahrräder sind ihre} & The bicycles are theirs\\\hline
        \bt{Das Buch ist mein} & The book is mine
    \end{tabular}
\end{example}
The forms for the `personal' and `possessive' pronoun classes are not distinguished in \bt{Deutsch}, but are clearly different. The roots for $\vm{P}[a,De,K\{M\}]$ are based upon the corresponding form in $\vm{P}[a,In,K](4,M)$, where $K$ represents the person being considered and $M$ its multiplicity. These roots  decline according to the negative article \bt{kein}. The polite variants $\vm{P}[a,De,2_p\{1\leftrightarrow\mathbb{M}, \mathbb{H}\}]$ require compulsory capitalization. An example for $\vm{P}[a,De,2_p\{\mathbb{M},0\}]$ is in \autoref{tab:denounT7}.
\begin{table}[bth]
    \centering
    \begin{tabular}{|c|c|c|c|}
        \hline\bt{Männ} & \bt{Weib} & \bt{Säch} & \bt{Mehr}\\\hline
        euer & eure & euer & eure\\\hline
        euren & eure & euer & eure\\\hline
        eurem & eurer & eurem & euren\\\hline
        eures & eurer & eures & eurer\\\hline
    \end{tabular}
    \caption{$2_p$ Plural possessive pronoun}
    \label{tab:denounT7}
\end{table}
Some clarifying examples:
\begin{example}
    \begin{tabular}{c|c}
        \bt{(Meine) Fahrräder haben keine Makel} & My bicycles have no faults\\\hline
        \bt{Sie hat sie unter (ihrer) Obhut} & She has her under her care\\\hline
        \bt{Alle lieben (Ihre) Vorträge} & All love his/her lectures
    \end{tabular}
\end{example}
Reflexive pronouns are encompassed under \bt{sich} and its forms as listed in \autoref{tab:denounT8}. These pronouns just imply self-referencing, and cannot be used with every verb. Their placement is always exactly after the conjugating verb or its illusion (in case it is exiled to the end of the sentence).
\begin{table}[bth]
    \centering
    \begin{tabular}{|c|c|c|c|c|c|c|}
        \hline $\Diamond$ & \multicolumn{2}{|c|}{$K=1_p$} & \multicolumn{2}{|c|}{$K=2_p$} & \multicolumn{2}{|c|}{$K=3_p$}\\\hline
        \bt{Akk} & mich & uns & dich & euch & sich & sich \\\hline
        \bt{Dat} & mir & uns & dir & euch & sich & sich \\\hline
    \end{tabular}
    \caption{Reflexive pronouns: $\vm{P}[r,De,K]$}
    \label{tab:denounT8}
\end{table}
\begin{example}
    \begin{tabular}{m{7cm}|m{7cm}}
        %\bt{Er stellt sich hinter mich} & He positioned himself behind me\\\hline
        \bt{Ich wasche mir die Hände}& I wash \ut{using myself the hands}\\\hline
        \bt{Wir treffen uns am Montag} & We \ut{are meeting us} on Monday\\\hline
        \bt{Ich bedanke mich fürs\{für + das\} Zuschauen} & I \ut{give the thanks within myself to} \ut{(those who saw)} for the viewing
    \end{tabular}
\end{example}
The last example  uses a special  construction used for fixed reflexiveness, something not found in English. It essentially means, `the thanks came from him' but it never mentions to whom it is given, implying it as obvious information. Here, the choice of the reflexive pronoun must additionally match  the fixed case code for the concerned usage. In this case, \bt{Ich bedanke mich} demands case $2$, so \bt{ich bedanke mir} is invalid. Occurrences are rare and the advice is to memorize these. Formal \bt{Sie} reflexives are  $\vm{P}[r,De,3_p]$ forms: these do not need capitalization.

Indefinite pronouns change forms according to corresponding article templates in both dependent and independent uses, but these changes are not consistently assigned. Some indefinite pronouns  can be used exclusively for objects, some exclusively for people, while most others fit for both. \autoref{tab:denounT9} documents various examples along with their followed  form variations. If any adjectives are used in the noun unit during such dependent use, they can be  declined in one of two ways:
\begin{enumerate}
    \item Assume the pronoun is the template word, and decline normally
    \item Use `without article' declinations
\end{enumerate}
The first rule merely assumes that the pronoun is structurally the same as the word, whose form template it follows. Both these rules are mutually exclusive, but have been assigned rather randomly. To mark pronouns which follow the second rule for adjective form variation, an asterisk (*) is added to their entries in \autoref{tab:denounT9}.
\begin{example}
    \begin{tabular}{c|c}
        [Aus (welchem) seltsamen Grund] benutzt du deine Sphäre nicht? & der\\\hline
        [Für (welche) seelische Flamme] möchtest du dich zerkleinern? & die\\\hline
        [(Manche) teuere Flaschen] sind es nicht wert & die \\\hline
        [(Solche) einzelne Menschen] sind meist Feiglinge & der
    \end{tabular}
\end{example}
The noun units and corresponding indefinite pronouns have been marked. It is also important to note that some of these pronouns may also directly be used as normal adjectives: \bt{Eine solche Sache, die wenigen Leute}\ldots clearly, the separation is fluid, and the only practically important thing is the applicable structure per use case.
\begin{table}[bth]
    \centering
    \begin{subtable}{\textwidth}
        \centering
        \begin{tabular}{|c|c|c|}
            \hline $\diamondsuit$ & \bt{Dependent} & \bt{Independent}\\\hline
            Alle & die (Mehrzahl)& die (Mehrzahl)\\\hline
            Einige* & die (Mehrzahl) & die (Mehrzahl)\\\hline
            Irgendeiner & ein/eine & der/die/das (Einzahl)\\\hline
            Jeder & der/die/das (Einzahl) & der/die/das (Einzahl)\\\hline
            Keiner & kein/keine & der/die/das \\\hline
            Mancher* & der/die/das & der/die/das\\\hline
            Solcher & der/die/das & der/die/das \\\hline
            Derartig* & Ohne Gesch.wort in \autoref{tab:denounT5} & Ohne Gesch.wort in \autoref{tab:denounT5}\\\hline
            Mehrere* & die (Mehrzahl) & die (Mehrzahl)\\\hline
        \end{tabular}
        \caption{Universally applicable}
    \end{subtable}
    \begin{subtable}{\textwidth}
        \centering
        \begin{tabular}{|c|c|c|}
        \hline $\clubsuit$ & \bt{Dependent} & \bt{Independent}\\\hline
        Nichts & ! (Ungültig) & - (Keine Formänderung)\\\hline
        Alles & ! & das (Einzahl)\\\hline
        Welcher & der/die/das & der/die/das \\\hline
        Irgendwelcher & der/die/das & der/die/das\\\hline
        Viel/Viele* & - & - \\\hline
        Wenig & ! & - \\\hline
        Etwas & ! & -\\\hline
        \end{tabular}
        \caption{Only inanimate objects}
    \end{subtable}
    \begin{subtable}{\textwidth}
        \centering
        \begin{tabular}{|c|c|c|}
            \hline $\spadesuit$ & \bt{Dependent} & \bt{Independent}\\\hline
            Man & ! & ein (Männ, kein Fall $6$)\\\hline
            Jemand & ! & ein (Männ, kein Fall $6$)\\\hline
            Irgendjemand & ! & ein (Männ, kein Fall $6$)\\\hline
            Niemand & ! & ein (Männ, kein Fall $6$)\\\hline
            Irgendwer & ! & der ($\vm{Q}_{M}(:,1)$) \\\hline 
        \end{tabular}
        \caption{Only people}
    \end{subtable}
    \caption{Indefinite pronouns: $\vm{P}[b,De/In]$}
    \label{tab:denounT9}
\end{table}

Some supplementary explanations to \autoref{tab:denounT9} will help digest it better:
\begin{enumerate}
    \item \bt{Alles/Alle} - \bt{Alles} represents a whole of something, and strictly refers to non-living entities, whereas \bt{Alle} has no such restriction and represents multiple wholes. \begin{example}
        \begin{tabular}{c|c}
            \bt{Alles war zerstört} & Everything was destroyed \\\hline
            \bt{Alle waren nass} & All were \ut{being} wet 
        \end{tabular}
    \end{example}
    \item \bt{Derartig} is an abbreviation of \textgerman{\enquote{\bt{von dieser Art}}} and declines like an adjective irrespective of whether it is used as a pronoun or adjective. The initial \bt{der} is not to be touched, only the endings vary.
    \item The prefix \bt{irgend-} (some) may be added to most of the W-words like `when, why, who\ldots' to create `someone, something, somebody\ldots' and so on. Using \bt{nirgend-} negates the corresponding word: `nobody, nowhere\ldots' \bt{Irgend-/Nirgend-} are merely  significance markers: using them on any word only decreases its perceived significance. Ex: \begin{example}
        \begin{tabular}{c|m{7cm}}
            \bt{Jemand ist angekommen} & Someone (Unknown) has arrived \\\hline
            \bt{Irgendjemand ist angekommen} & Someone (Unknown and irrelevant) has arrived
        \end{tabular}
    \end{example}
    These \bt{irgend-} compounds never decline. The only ones who do are documented in \autoref{tab:denounT9}. Commonly appended pronouns are: \bt{etwas, wo, wie, was, wann}.
    \item The diligent reader must have realized that the categories of `pronoun, adjective, article' are not completely independent. As long as the meaning supports it, any word may be a member of multiple  classes, where it will follow the respective structure. Ex: \begin{example}
        \begin{tabular}{c|c|c|c}
            \bt{Base} & \bt{Pronoun} & \bt{Adjective} & \bt{Article}\\\hline
            ein & einer & ein & ein \\\hline
            viel & viel & viele & ! 
        \end{tabular}
    \end{example}
    \bt{Viel/Viele} are the same as  much/many, and only mark countability. Since no emphasis is marked, article usage becomes invalid.
    \item \bt{Wer} requires hook word forms which will be introduced in \autoref{sec:destce}.
\end{enumerate}

\autoref{tab:denounT10} details the  demonstrative pronouns along with their corresponding form variations.
\begin{table}[bth]
    \centering
    \begin{tabular}{|c|c|c|c|}
        \hline $\heartsuit$ & \bt{Dependent} & \bt{Independent} & \ut{Meaning}\\\hline
        Dies & der/die/das & kein/keine & \textit{This} \\\hline
        Das & ! & - & \textit{This/That}\\\hline
        \textcolor{magenta}{Derselbe} & \autoref{tab:denounT5}, \autoref{tab:denounT4} & \autoref{tab:denounT5}, \autoref{tab:denounT4} & \textit{the same}\\\hline
    \end{tabular}
    \caption{Demonstrative pronouns: $\vm{P}[c,De/In]$}
    \label{tab:denounT10}
\end{table}
Both \bt{das} and \bt{dies} have the same meaning, but with different emphases. \bt{Das} is rather neutral and general, while \bt{dies} is formal and much more concrete. \textcolor{magenta}{Derselbe} is a fusion of the definite article and the adjective \bt{selb}, meaning `same'. This word follows all noun unit rules for articles and adjectives, and implies that something is exactly something else, which should not be confused with \bt{das gleiche}, which implies that something is exactly equal to something else. \bt{Derselbe} implies that the two entities being considered are one and the same, not merely equal.
\begin{example}
    \begin{tabular}{m{7.5cm}|m{7.5cm}}
        \bt{Grober Kies liegt auf diesem Strand} & Coarse gravel lies on this beach\\\hline
        \bt{Das kann unseren Einsatz gefährden} & This could \ut{our mission endanger}\\\hline
        \bt{Zwillinge entstehen aus demselben Mutterleib, und genießen oft die gleiche Gunst} & Twins emerge out of the same womb, and \ut{enjoy often} the same favour
    \end{tabular}
\end{example}
An exact analogue to a far-away marker like `that' does not exist.  Adverbs like \bt{drüben, dort, da}  concretely mark distance, and are used with \bt{das} to get `that'. Ex: \bt{das dort, dort, drüben, das da, da drüben\ldots}

\subsection[Joining]{Joining (Wortverkettung)}
I hinted earlier that \bt{Deutsch} uses a lot of concatenated words that encode meanings via processes. Compared to the sheer amount of such compounds, the rules concerning them are surprisingly compact:
\begin{enumerate}
    \item Add \bt{-s} to the end of the first word and join
    \item Add \bt{-en/-n} to the end of the first word and join
    \item If the first word is a verb, eliminate its infinitive \bt{-en} and join
    \item Join
\end{enumerate}
These four rules are mutually exclusive (obviously), but there is no real pattern as to which applies where. A lot of \bt{Deutsch} has old relics still in use,  even though the language evolved and nobody uses the old rules anymore. New words are formed using new rules, and the ones which used the old rules remain as exceptions\ldots of which there are many. Giving some examples is all I can do for this section:
\begin{example}
    \begin{tabular}{c|c|c}
        Anmeldung(s)/zeit/raum & Flug/zeug & Gär(-en)/kessel\\\hline
        Tag(es)/licht & Schnecke(n)/haus & Wiedergabe/geschwindigkeit\\\hline
        Krank(en)/wagen & Tür/störung & Fahr(-en)/strecke
    \end{tabular}
\end{example}
You will always join the singular of the first word. The compound can then be treated like a regular noun. This compound noun takes the gender and word forms of the last noun in the chain. Oh, and there is technically no limit to how many of those you can smash together as long as it makes sense. This convenience is above the mental capacities of some people, which is why German gets insulted by the inept for its long words, complicated nature and aggressiveness. It is also likely because of these long words that Germans are habituated to speaking fast, else sentences would take a while to complete.

Mostly, the first word has cases $6$ or $4$ if the compound were to be expanded. Keep in mind  that the associated case is not necessarily unique, since the expansion is not unique. It just has to make sense. This feature `cleans' the language by introducing roughness, which is why English speakers  struggle initially.
\begin{example}
    \begin{tabular}{m{7cm}|m{7cm}}
        \bt{Anmeldungszeitraum}  & Registration time period\\\hline
        \bt{Flugzeug} &  Airplane (Flight thing)\\\hline
        \bt{Wiedergabegeschwindigkeit} &  Playback speed
    \end{tabular}
\end{example}
If the meaning allows it, adjectives/nouns  can also be pre-joined to other adjectives. In such cases, \ut{only} joining rule `4' is allowed:
\begin{example}
    nagelneu, blitzsauber, uralt, bildschön, allerschönst, größtmöglich\ldots
\end{example}
\subsection[Numbers]{Numbers (Zahlen)}
Before starting this subsection, it must be declared that the gender of German numbers, if at all required is always \bt{sächlich}. Let us begin with pure numbers. These  never show form variation, even when placed in noun units (except \bt{eins}, which  becomes the indefinite article \bt{ein}). Such numbers are used to portray time. If written as digits, \bt{Deutsch} uses a comma as a decimal point.
\begin{example}
eins, drei, zwanzig, zweiundzwanzig, achtundvierzig, einundfünfzig, elf\ldots
\end{example}
To count ranks, there are two rules:
\begin{enumerate}
    \item $\{1,2,\ldots, 19\}$: Add \bt{-t} at the end of the pure number. $1$ and $3$ are the only exceptions with counters \bt{erst,  dritt} respectively.
    \item $\{20,21,\ldots\}$: Add \bt{-st} at the end of the pure number.
\end{enumerate}
These rank counter root forms behave like adjectives, and are subject to their rules. It is possible to say things like `as a couple, as a quintuple' by pairing \bt{zu} and the respective rank root. This usage does not undergo form variation.
\begin{example}
    zu zweit, zu dritt, zu fünft, zu sechst\ldots
\end{example}

To define  multiplicities,  the suffix \bt{-fach} (times) is added after  the factor $f$ as defined in \autoref{eq:factor}. $[f]$ is represented by pure numbers. Fractional bases use the suffix \bt{-el} at the end of the corresponding rank counter, with their multiplicity (as a pure number)  appended at the beginning to make $\{f\}$. $f$ is created by joining both of these  in the order $([f], \{f\})$ with a  dash between the components to ensure readability:
\begin{example}
    \begin{tabular}{c|c}
    \bt{Integral} & (einzeln), zweifach, dreifach, vierfach, einundfünfzigfach\ldots\\\hline
    \bt{Fractional} & (einhalb \{$\frac{1}{2}$\}), dreiviertel ($\frac{3}{4}$), neundreiundachtzigstel ($\frac{9}{83}$)\ldots\\\hline
    \bt{Mixed} & zwei-eindrittel-fach ($2\frac{1}{3}$), fünfundfünfzig-dreiviertel-fach ($55\frac{3}{4}$)\ldots
    \end{tabular}
\end{example}
All of these behave like adjectives. `Single' can be written as \bt{einfach} but is usually represented by \bt{einzeln}, while \bt{einfach} means `simple'. 

Verbs that represent the action of multiplying by a factor have the pattern:
\begin{syntax}
    \begin{tabular}{c|c}
    Multiplication: &\bt{ver +} Factor \bt{+ fach + en}\\\hline
    Division: & Pure number \bt{+ teilen}
    \end{tabular}
\end{syntax}
This then gives rise to verbs like: \bt{verzweifachen/verdoppeln,  verzweiundzwanzigfachen, verzwei-eindrittel-fachen, vervierelftel-fachen, (halbieren)\ldots}. Be mindful that fractional components always need \textit{CoN} separation (-). Pure number division verbs denote the division of an object into its argument number of parts: \bt{zweiteilen, vierteilen, zweiundzwanzigteilen\ldots}. These verb classes are regular and use the helping verb \bt{haben}.
\section[Verb]{Verb (Tätigkeitswort)} \label{sec:deverb}
An infinitive (dictionary form/\bt{Grundform}) has the ending \bt{-en/n}.  For every verb $V\in \mathcal{S}$, all structural rules are exactly the same. Helping verbs ($\mathcal{X}$) are used to specify tense, sentence mode, or  both. Members of these sets are:
\begin{setdef}
    \begin{tabular}{m{14cm}}
        $\mathcal{S}$ : \{Müssen, Können, Dürfen, Mögen, Sollen, Wollen, Lassen\}\\\hline
        $\mathcal{X}$ : \{Haben, Sein, Werden\}
    \end{tabular}
\end{setdef}

Before explaining the rules pertaining to $\mathcal{Z}$, it is useful to consider a small fact about the VC: The DC is always exiled to the end of any sentence, while the CV stays at second place. This means that all verbs that take arguments will have them  thrown to the very end irrespective of use case. The set $\mathcal{Z}$ is labelled in \bt{Deutsch} as \bt{trennbar} (separable):
\begin{example}
    \begin{tabular}{c c|c c}
        \textit{Verb} & \textit{Meaning} & \textit{Verb} & \textit{Meaning}\\\hline
        \bt{auf/geben}& to give up & \bt{unter/legen} & to place under \\\hline
        \bt{an/ziehen} & to pull towards (attract) & \bt{hin/schreiben}& to write down\\\hline
        \bt{auf/hören}& to hear up (stop) & \bt{an/lernen} & to learn towards (train)
    \end{tabular}
\end{example}
As seen in the last  row, some prepositionally supported verbs mean something  different than the sum of their parts. This becomes especially annoying due to the counterintuitive placement of the DC at the end, which is why using and interpreting verbs in $\mathcal{Z}$ takes some getting used to. The set $\mathcal{Y}$ is labelled as \bt{untrennbar} (inseparable), and behave like typical verbs with prefixes. These can themselves have prepositional arguments: \bt{an/vertrauen} (to entrust), \bt{an/erkennen} (to recognise)\ldots. Sometimes, they also have arguments that behave like prepositions but aren't: \bt{überein/stimmen} (to match), \bt{kennen/lernen} (to get to know someone)\ldots

A pronunciation tip: The preposition (no matter where it is) is stressed in \bt{trennbar} verbs whereas the main stem and not the prefix is stressed in \bt{untrennbar} verbs. If there are multiple prefixes (rare), then the stress alternates like \bt{\textit{miss}-ver-\textit{stehen}}, with  italic text showing the stressed sections.

For advanced learners, some verbs belong to both $\mathcal{Z}, \mathcal{Y}$, and both use cases mean different things. This is not completely random: $\mathcal{Z}$  is literal, whereas $\mathcal{Y}$ is rather symbolic:
\begin{example}
    \begin{tabular}{c|c|c}
$\mathcal{Z}$ & \bt{Er schaute die Unterlagen durch} &  He was looking through the documents\\\hline
$\mathcal{Y}$ & \bt{Sie hat seine Pläne durchschaut} & She has \ut{his plans} seen through 
    \end{tabular}
\end{example}
Such verbs have some specific associated prepositions: \bt{durch, über, unter, um}. So you can first suspect foul play when you see any of these, while the general prepositional argument only creates a member of $\mathcal{Z}$. Additionally, among such verbs with dual citizenship, the helping verb for building the \bt{Perfekt} form of $\mathcal{Y}$ is always \bt{haben}.

Like English, \bt{Deutsch} does not differentiate between object and non-object verbs, except by context: \bt{erschrecken} (to frighten/to be frightened). The object for any object verb is a case $2$ unit \textit{without} an associated preposition. This distinction is important, since case $2$ can otherwise appear with a preposition even if a non-object verb is used (\bt{Akk/Dat} - \autoref{sec:despco}).
\subsection[Tense]{Tense (Handlungsart und Zeitstufe)}
\autoref{tab:deverbT1} represents the German aspect-time matrix. Any forms for the polite \bt{Sie} can be obtained by letting the CV stay in its infinitive form, irrespective of which VC is being considered. Hence, their explicit description can be safely ignored to declutter the matrix.
\begin{table}[bth]
    \centering
    \begin{tabular}{|c|c|c|c|c|}
        \hline $\Sigma $ & Past & Present & Future & Neutral\\\hline
        Ongoing & \bt{Präteritum} & \bt{Präsens} & \bt{Futur 1} & \bt{zu + Grundform}\\\hline
        Completed & \bt{Plusquamperfekt} & \bt{Perfekt} & \bt{Futur 2} & \bt{MdV + zu + $\mathcal{X}$}\\\hline
    \end{tabular}
    \caption{Aspect-Time matrix: $\vm{T}$}
    \label{tab:deverbT1}
\end{table}

Pure terms for the tenses in \autoref{tab:deverbT1} have not been used  because they are unoriginal and too long. The base cell is $\vm{T}(1,2)$, while the \bt{Grundform} refers to the verb infinitive. 

\textcolor{magenta}{\bt{First, aspect}}  constructions for verbs $V\notin \mathcal{S}$ and $\vm{T}(:,2)$ are considered:
\subsubsection{$\vm{T}(1,2)$: Present Ongoing (Präsens)}
Denotes an action that is occurring at the time of expression. This is the base cell, with  form variation as in \autoref{tab:deverbT2}.
\begin{table}[bth]
    \centering
    \begin{tabular}{|c|c|c|}
        \hline $\zeta$ & $1$ & $\mathbb{M}$\\\hline
        $1_p$ & VR + \bt{e} & \bt{Grundform}\\\hline
        $2_p$ & VR + \bt{st} & VR + \bt{t}\\\hline
        $3_p$ & VR + \bt{t} & \bt{Grundform}\\\hline
    \end{tabular}
    \caption{Aspect: Ongoing, $\vm{S}[\vm{T}(1,2)]$}
    \label{tab:deverbT2}
\end{table}

All conjugations are based upon the verb root (VR), and some VRs show exceptional configurations exactly for $\vm{S}[\vm{T}(1,2)](2_p,1,0)$ and $\vm{S}[\vm{T}(1,2)](3_p,1)$. Since there is only one verb in the chain, it is also the CV, while the DC is a null set. An example for both regular and irregular form variation:
\begin{example}
    \begin{tabular}{c c}
    hören - \bt{to hear}&
    \begin{tabular}{c|c|c}
        $\pi$ & $1$ & $\mathbb{M}$\\\hline
        $1_p$ & höre & hören \\\hline
        $2_p$ & hörst & hört \\\hline
        $3_p$ & hört & hören
    \end{tabular}\\\hline
    erlöschen - \bt{to be extinguished}&
    \begin{tabular}{c|c|c}
        $\pi$ & $1$ & $\mathbb{M}$\\\hline
        $1_p$ & erlösche & erlöschen \\\hline
        $2_p$ & erlischst & erlöscht \\\hline
        $3_p$ & erlischt & erlöschen
    \end{tabular}
\end{tabular}
\end{example}
For verbs whose \bt{Grundform} ends with  \bt{-eln}, the  \bt{e} is eliminated from the VR in $\vm{S}[\vm{T}(1,2)](1_p,1)$, merely for pronunciation purposes. Ex: \bt{krabbeln/krabble}
\begin{syntax}
    CV Conjugation pattern (CVC): $\vm{T}(1,2)$
\end{syntax}
Example sentences:
\begin{example}
    \begin{tabular}{c|c}
        \bt{Ich vergesse mich} & I am forgetting myself\\\hline
        \bt{Er reinigt sein Zimmer} & He is cleaning his room
    \end{tabular}
\end{example}

\subsubsection{$\vm{T}(2,2)$: Present Completed (Perfekt)}
Denotes a completed action at the time of expression. Note that it does not matter how recently the action was completed, just that it is complete at the considered  moment of expression. Hence, \bt{Perfekt} coupled with adverbs is often  used to express things that happened in the past, although it is not technically a past form. The way this tense is built is twofold:
\begin{enumerate}
    \item Choose the correct  helping verb for the CV from $\mathcal{X}$: \bt{haben/sein}
    \item Convert the CV into its \bt{MdV} form
    \item The \bt{MdV} becomes appended to the end of the DC, while the helping verb becomes the new CV
\end{enumerate}
Depending on the CV, either of \bt{haben/sein} can be chosen as the helping verb according to two universal rules:
\begin{syntax}
    For any verb $V$, if
    \begin{itemize}
        \item $V$ does not take case $2$ nouns ($V$ is a non-object verb) \ut{\textit{AND}}
        \item The case $1$ noun used with $V$ [effects a change of state or non-oscillatory motion on itself / undergoes a change of state or non-oscillatory motion]
    \end{itemize}
    then $V$ uses helping verb \bt{sein}. Otherwise, $V$  uses \bt{haben}.
\end{syntax}
Exactly two exceptions exist which use \bt{sein}: \bt{bleiben, sein}.  Around 80\% of all verbs do not satisfy both of the aforementioned conditions simultaneously, so if you find yourself  unsure of what to use, default to \bt{haben}.

The \bt{Mittelwort der Vergangenheit (MdV)}, which  grammar textbooks simply call \bt{Partizip2} is the analogue of P2 (gone, looked\ldots) in English. The \bt{MdV} is formed by either of three rules, also depending on the irregularity of the verb. For irregular verbs, the VR undergoes transformation.
\begin{enumerate}
    \item \bt{ge} + VR + \bt{t}
    \item \bt{ge + Grundform}
    \item \textcolor{magenta}{VR + \bt{t}}
\end{enumerate}
\textcolor{magenta}{Only} verbs that end in \bt{-ieren} like: \bt{stornieren, fungieren, studieren, fundieren} exclusively follow the third pattern. These are mostly germanised verbs and although I personally dislike them, that is no reason to leave them out.

All other verbs follow either of the first two patterns. Irregular verbs use their \bt{Präteritum} forms, or  close variants to them as the VR. They do have certain repeating patterns, but it is not possible to describe these, and  are best learned through observation. Any verbs in $\mathcal{Y}$ do not take the initial \bt{ge-}, since their own prefix overrides its addition. Exceptions are incredibly rare: \bt{gekennzeichnet, gehandhabt}. Verbs in $\mathcal{Z}$ behave as usual, and the argument simply attaches itself before the \bt{MdV} form. Some verb/\bt{MdV} pairs:
\begin{example}
    haben/gehabt, sein/gewesen, lassen/gelassen, erstarren/erstarrt, verlassen/verlassen, entscheiden/entschieden, geschehen/geschehen, sterben/gestorben, aufgeben/aufgegeben, kennenlernen/kennengelernt, übereinstimmen/übereingestimmt, emporwinden/emporgewunden
\end{example}
If a verb can be used as both an object and non-object verb, sometimes the \bt{MdV} and \bt{Präteritum} forms corresponding to both variants are different:
\begin{example}
    \begin{tabular}{c|c|c}
        \bt{Verb} & \bt{Object:} MdV, $\vm{T}(1,1)$ \bt{root} & \bt{Non-object:} MdV, $\vm{T}(1,1)$ \bt{root}\\\hline
        erschrecken & erschreckt, erschreckt & erschrocken, erschrak \\\hline
        schmelzen & geschmelzt, schmelzt & geschmolzen, schmolz
    \end{tabular}
\end{example}
\begin{syntax}
    CV Conjugation pattern (CVC): $\vm{T}(1,2)$
\end{syntax}
Example sentences:
\begin{example}
    \begin{tabular}{c|c}
        \bt{Ich habe sie gestern gesehen} & I have  her yesterday \ut{seen}\\\hline
        \bt{Meine Mutter hat mein Fahrrad verkauft} & My mother has my bicycle \ut{sold}
    \end{tabular}
\end{example}
\textcolor{magenta}{\bt{Next, time}}  constructions for verbs $V\notin \mathcal{S}$ and $\vm{T}(1,:)$ are considered:
\subsubsection{$\vm{T}(1,1)$: Past Ongoing (Präteritum)}
Used for narrations. To get the \bt{Präteritum} root, add \bt{-t} as a suffix to the VR. The VR changes if the verb is irregular. Leaving aside the irregular formations of the roots, there are otherwise no conjugation irregularities. The pattern follows \autoref{tab:deverbT2}, but with some differences:
\begin{enumerate}
    \item $(3_p,1)$ follows $(1_p,1)$
    \item If the verb is irregular, $(1_p,1)=$ `\bt{Präteritum} root'
    \item Sometimes, an additional \bt{e} is added at the end of the \bt{Präteritum} root in $(2_p,:,0)$ for ease of pronunciation
\end{enumerate}
Two examples are given below to illustrate:
\begin{example}
    \begin{tabular}{c c}
    kaufen - \bt{to buy}&
    \begin{tabular}{c|c|c}
        $\pi$ & $1$ & $\mathbb{M}$\\\hline
        $1_p$ & kaufte & kauften \\\hline
        $2_p$ & kauftest & kauftet \\\hline
        $3_p$ & kaufte & kauften
    \end{tabular}\\\hline
    sterben - \bt{to die}&
    \begin{tabular}{c|c|c}
        $\pi$ & $1$ & $\mathbb{M}$\\\hline
        $1_p$ & starb & starben \\\hline
        $2_p$ & starbst & starbt \\\hline
        $3_p$ & starb & starben
    \end{tabular}
\end{tabular}
\end{example}
Common irregular patterns used to make \bt{Präteritum} roots are:
\begin{setdef}
gehen/ging, sterben/starb, sehen/sah, frieren/fror, halten/hielt, werden/wurde, reißen/riss, entscheiden/entschied, lassen/ließ
\end{setdef}
\begin{syntax}
    CV Conjugation pattern (CVC): $\vm{T}(1,1)$
\end{syntax}
Example sentences:
\begin{example}
    \begin{tabular}{m{5.5cm}|c}
        \bt{Ich ging nach dem Supermarkt} & I was going to the supermarket\\\hline
        \bt{Meine Hoffnungen starben am(an+dem) Dienstag} & My hopes were dying on Tuesday
    \end{tabular}
\end{example}
\subsubsection{$\vm{T}(1,3)$: Future Ongoing (Futur 1)}
Uses the helping verb \bt{werden} for the purpose. The formation proceeds as follows:
\begin{enumerate}
    \item Append the CV in its \bt{Grundform} to the end of the DC
    \item The new CV is \bt{werden}, which follows \autoref{tab:deverbT2}
\end{enumerate}
\begin{syntax}
    CV Conjugation pattern (CVC): $\vm{T}(1,2)$
\end{syntax}
Example sentences:
\begin{example}
    \begin{tabular}{m{6cm}|c}
        \bt{Ich werde nach dem Supermarkt gehen} & I will be  to the supermarket \ut{going}\\\hline
        \bt{Du wirst bald sterben} & You will be \ut{soon} dying
    \end{tabular}
\end{example}
With that, all the cells in $\vm{T}$ can be successfully derived, even for varying sentence modes. For the sake of completeness:
\begin{example}
    \begin{tabular}{m{5cm}|c|m{5cm}}
    Ich werde nach dem Supermarkt gegangen sein & Futur 2 & \bt{I will to the supermarket \ut{gone have}}\\\hline
    Ich war nach dem Supermarkt gegangen & Plusquamperfekt & \bt{I had to the supermarket \ut{gone}}
    \end{tabular}
\end{example}
\subsubsection{Special verbs (Sondertätigkeitswörter)}
For verbs $V\in \mathcal{S}$, $V$ is the CV while the main verb in its \bt{Grundform} is the DC for the base cell: $\vm{T}(1,2)$. All transformations work as described earlier, but there are two special rules:
\begin{enumerate}
    \item If $V$ is sent to the DC, it \ut{MUST} stay at its very end in its \bt{Grundform}
    \item If any other $G\in\mathcal{X}$ is sent to the DC while $V$ is already in it, $G$ will not interact with the main verb or $V$, and will permeate as far as possible to the front of the DC
\end{enumerate}
In practice, the second rule is only applicable for $\vm{T}(2,3)$. The exact nature of interactions to be expected will be clear in \autoref{sec:destce}. Some examples to demonstrate both rules:
\begin{example}
    \begin{tabular}{m{6.5cm}|m{5.5cm}|c}
    Wir werden lange haben warten [müssen] & \bt{We will long \ut{have waited} \ut{[have to]}} & Futur 2\\\hline
    Ich habe heute gehen [sollen] & \bt{I \ut{have} today \ut{gone} \ut{[should]}} & Perfekt\\\hline
    Du darfst die Gesellschaft nicht töten & \bt{You are allowed to society \ut{not kill}} & Präsens
    \end{tabular}
\end{example}
The last thing to mention about $\mathcal{S}$ is that these verbs can be chained. One only needs to create the corresponding base cell according to which verb is used first, after which the above methods can be used directly:
\begin{example}
    \begin{tabular}{m{15cm}}
Let $S_1, S_2, S_3\in\mathcal{S}$, and one desires to use them with the main verb $M_V$ in the order $S_2, S_3, S_1$, with the left most component being used first. The corresponding base cell will be:\\
CV: $S_1$, DC: $M_V + S_2 + S_3$\\
\bt{Perfekt:} CV = `haben', DC = $M_V + S_2 + S_3 + S_1$\\ \bt{Futur 2:} CV = `werden', DC = $\texttt{haben} + M_V + S_2 + S_3 + S_1$
    \end{tabular}
\end{example}
Practically, I do not believe it is possible to chain more than two meaningfully, but a complicated case is always better for illustration. Also, since all verbs in $\mathcal{S}$ use the helping verb `haben' for building their \bt{Perfekt} aspect, a direct substitution has been made in the above example.

\subsection[Adverb]{Adverb (Umstandswort)}
There are some important formation patterns based on the word ending. Since there is a significant overlap between the set of all adjectives and adverbs,  the following patterns are to be assumed applicable solely for adverbs, unless otherwise specified. Note that this list is not comprehensive but  sufficient:
\begin{enumerate}
    \item Unlike English, \bt{Deutsch} does not differentiate between `clear' and `clearly'. Hence, most adjectives also double as adverbs: 
    \begin{example}
    \bt{herrlich} (magnificently), \bt{seltsam} (strangely), \bt{begrenzbar} (terminably)\ldots
    \end{example}
    \item \bt{-s}: Normally have something to do with placement in time/space, and made from nouns: 
    \begin{example}
    \bt{morgens} (in the morning), \bt{abends} (in the evening), \bt{rechts} (to the right)\ldots
    \end{example}
    \item \bt{Eigenschaftswort + er + maßen}: Shows  something  in accordance to the associated adjective: 
    \begin{example}
    \bt{erwiesenermaßen} (as is proved), \bt{gewissermaßen} (to a certain extent)\ldots
    \end{example}
    \item \bt{-weise}: Denotes the manner in which  something proceeds. If the root word is an adjective, the construction follows the previous rule: 
    \begin{example}
    \bt{normalerweise} (normally), \bt{ähnlicherweise} (similarly), \bt{zeitweise} (from time to time)\ldots
    \end{example}
\end{enumerate}
Like with verbs, only  adverb units can be used in sentences. 

\subsection[Nominalization]{Nominalization (Nennwortbildung)}
\subsubsection{Contraction (Verkürzung)}
To contract a noun, take the normal noun unit, eliminate the noun, and then capitalize the adjective. The grammatical gender corresponds to the actual gender of the entity, or if the entity does not have a natural gender, its grammatical gender. The only exceptions to this rule are instruments like `cutter, accelerator' which are always \bt{männlich}, and follow the corresponding \bt{MdG} type contraction rules.

The contracted adjective always declines as if the omitted noun still exists.
\begin{example}
    \begin{tabular}{c c}
    Pure adjective & \begin{tabular}{c|c}
        \textit{Expanded} & \textit{Contracted}\\\hline\bt{Das hübsche Mädchen} & \bt{Die Hübsche}\\\hline \bt{Gefährliche Frauen} & \bt{Gefährliche} \end{tabular}\\\hline
    \bt{MdG} & \begin{tabular}{c|c}
        \textit{Expanded} & \textit{Contracted}\\\hline\bt{Im folgenden Teil} & \bt{Im Folgenden} \\\hline \bt{Zu der als Anwältin arbeitenden Frau} & \bt{Zu der Anwältin}\\\end{tabular}\\\hline
    \bt{MdV} & \begin{tabular}{c|c}
        \textit{Expanded} & \textit{Contracted}\\\hline \bt{Die gefangenen Bevölkerungen} & \bt{Die Gefangenen}\\\hline\bt{Die ungewollten Leute}& \bt{Die Ungewollten}\end{tabular}
    \end{tabular}
\end{example}
Some explanations are in order. 
\begin{itemize}
    \item \bt{MdG (Mittelwort der Gegenwart)} is the German counterpart to the P1 form
    \item \bt{MdG} type contractions representing male professions  use the suffixes \bt{-er/-er} on the verb root for \bt{Ein/Mehrzahl}. If there is no associated verb, then the pattern \bt{Sturz/Stürze} is followed instead. Female professions always use suffix \bt{-in} on the corresponding male profession for building the \bt{Einzahl}. The \bt{Mehrzahl} is built by \bt{-innen}. Sometimes, they also have an extra \bt{Umlaut}: \bt{Arzt/Ärzte, Ärztin/Ärztinnen}; \bt{Fahrer/Fahrer, Fahrerin/Fahrerinnen}
    \item If there are multiple adjectives in a noun unit, only one of them can be contracted according to the above rule, while the other behaves as normal
    \item Anything can be contracted  as long as the sentence makes sense in context
\end{itemize}
Nationalities are a special case for this rule. Even though they should technically be pure adjective type contractions, \bt{MdG} type nouns are used instead:
\begin{example}
    \begin{tabular}{c|c}
        \textit{Male/Female} & \textit{Nationality}\\\hline
        \bt{Russe/Russin} & Russian\\\hline
        \bt{Inder/Inderin} & Indian \\\hline
        \bt{Franzose/Französin} & French
    \end{tabular}
\end{example}
The only country for which this works normally is Germany: \bt{Deutscher/Deutsche}
\subsubsection{Conceptualization (Vorstellungsbildung)}
A collection of all entities that possess a certain quality can be described by using the contraction rule on a pure adjective, but keeping the grammatical gender as \bt{sächlich}. No associated noun exists in this case, but form changes still assume that the noun unit exists as usual:
\begin{example}
    \begin{tabular}{c|m{3.5cm}|m{7cm}}
        \bt{Das Deutsche} & German (refers to the language) & \bt{Im Deutschen gibt es einen Makel}\\\hline
        \bt{Das Kluge} & The clever & \bt{Auf den Klugen kann man immer vertrauen}\\\hline \bt{Das Böse}& The Evil & \bt{Das Böse hat gewonnen}
    \end{tabular}
\end{example}
Suffixes  \bt{-heit/-keit} or a form change with \bt{Umlaut + e} are the patterns used for creating concept nouns. The suffixes \bt{-heit/-keit} can also be added to VRs and \bt{MdV}s, and the consequently formed nouns behave similarly. All of them  are always grammatically \bt{weiblich}:
\begin{example}
    Güte, Stärke, Schönheit, Sauberkeit, Trägheit, Verrücktheit, Besessenheit, Kleinigkeit, Verträglichkeit\ldots
\end{example}
Making process nouns from verbs is as follows:
\begin{enumerate}
    \item Capitalize the \bt{Grundform} and assume the gender is \bt{sächlich}
    \item Yeah. The entirety of $\vm{C}$ for these nouns except for $\vm{C}(4,1)$ is exactly the same
\end{enumerate}
The only reason $\vm{C}(4,1)$ even changes is due to \bt{Genetiv} rules. Some examples:
\begin{example}
    Denken, Fahren, Gehen, Zusammenleben, Berücksichtigen, Töten\ldots
\end{example}
State nouns  are mostly grammatically \bt{weiblich}, and are mostly formed by adding a suffix to the VR:
\begin{example}
    \begin{tabular}{c|c|m{9cm}}
        \textit{Suffix} & \textit{Frequency} & \textit{Nouns}\\\hline
        \bt{-ung} & Extreme & \bt{Bildung, Verwaltung, Einstellung, Zahlung, Fahndung, Ursprung}\\\hline
        \bt{-e} & Low & \bt{Sperre, Reise, Runde}\\\hline
        \bt{-} & Average & \bt{Bild, Urteil, Teil, Regel, Anruf}\\\hline
        \bt{-nis} & Low & \bt{Versäumnis, Erlaubnis, Erlebnis}\\\hline
        VR transform & Low & \bt{Verschluss, Tod}
    \end{tabular}
\end{example}
Lastly, there are only two noun suffixes worth considering:
\begin{enumerate}
    \item \bt{-chen}: Reduces scale. Gender \bt{sächlich}: 
    \begin{example}
        \begin{tabular}{c|c}
            \bt{Teilchen} & Particle[s]\\\hline
            \bt{Früchtchen}& Small fruit[s]\\\hline \bt{Brötchen} & Small bread[s]
        \end{tabular}
    \end{example}
    \item \bt{-schaft}: A collective/community; a state/condition; a property. This suffix can sometimes also apply to adjectives or adverbs. Gender \bt{weiblich}:
    \begin{example}
        \begin{tabular}{c|c}
            \bt{Gemeinschaft[en]} & Community[ies] \\\hline \bt{Herrschaft[en]} & Dominion[s] (The state of domination)\\\hline \bt{Freundschaft[en]} & Friendship[s] \\\hline \bt{Wissenschaft[en]} & Science[s] (Collection of knowledge)\\\hline \bt{Landschaft[en]} & Landscape[s]
        \end{tabular}
    \end{example}
\end{enumerate}
Of these, \bt{-schaft} is much more common and not as restricted in use as \bt{-chen}. 

Fun fact: All words that end in \bt{chen} are \bt{sächlich}, leading to \bt{Mädchen} (girl) being  grammatically neutral.
\subsection[Prefix]{Prefix (Vorsilbe)}
\autoref{tab:deverbT3} has the most common prefixes used to create \bt{untrennbar} verbs with their usual meanings. Though useful for a formal entry, there are many exceptions with regard to meaning, so be mindful.
\begin{table}[bth]
    \centering
    \begin{tabular}{|c|m{4.5cm}|c|}
        \hline\textit{Prefix} & \textit{Meaning}& \textit{Examples}\\\hline
        \bt{be} & Emphasizes the result of an action. This prefix always creates object verbs, and allows no prepositions in the case $2$ noun unit & \bt{beängstigen, beantworten, bereuen}\\\hline
        \bt{ent} & Take something away & \bt{entfernen, enteignen, entschlüsseln}\\\hline
        \bt{er} & Result of an action & \bt{erlassen, erhöhen, erachten}\\\hline
        \bt{ge} & No specific meaning & \bt{geschehen, gelingen, gestehen}\\\hline
        \bt{emp} & No specific meaning & \bt{empfangen, empören, empfinden}\\\hline
        \bt{miss} & Opposite / Failure & \bt{misslingen, missachten, missglücken}\\\hline
        \bt{wider} & Against & \bt{widersprechen, widerstehen, widerfahren}\\\hline
        \bt{zer} & Destruction  & \bt{zerstören, zerreißen, zerschlagen}\\\hline
        \bt{ver} & Disappear,  make a mistake, no specific meaning, destroy,  multiply, opposite & \bt{verwenden, verachten, vermählen}\\\hline
    \end{tabular}
    \caption{Prefix table}
    \label{tab:deverbT3}
\end{table}

Of these prefixes, \bt{ver} is the most versatile, since it can form many specific verbs using adjectives or adverbs. The pattern is:
\begin{syntax}
    \bt{ver + Eigenschaftswort/Umstandswort + en}
\end{syntax}
They are so versatile, that I suspect it is possible to write entire stories using only this verb variant:
\begin{example}
    vernachlässigen, verabscheuen, verbergen, verlegen, verkaufen, verderben, versiegeln, verunreinigen, verschmutzen, verschmelzen, verhören, verschmachten, verzieren, verdanken, verlesen\ldots
\end{example}
If anyone ever manages to succeed in this challenge, I would very much like to read the resulting abomination.
\section[Preposition]{Preposition (Verhältniswort)} \label{sec:deprep}
The versatility \bt{Deutsch} offers makes it impossible to list out all  prepositions, so only the main concepts along with the most frequently used candidates will be touched upon. 

First, we look at \bt{her/hin} correspondence:
\begin{enumerate}
    \item \bt{her}: Movement towards the speaker
    \item \bt{hin}: Movement away from the speaker
\end{enumerate}
These prefixes can be slammed upon prepositions to create new ones with further specification:
\begin{example}
    \begin{tabular}{c c|c c}
        \bt{Heraus} & Out towards me & \bt{Hinaus} & Out away from me (beyond)\\
        \bt{Herein} & In towards me & \bt{Hinein} & In away from me
    \end{tabular}
\end{example}
The word \bt{herein} is consequently used for inviting someone into the room you are currently in. This difference is subtle, and rarely captured in a translator. Since these are also prepositions, they can be coupled with \bt{trennbar} verbs:
\begin{example}
    \bt{herauf/beschwören}: To conjure up (towards me)
\end{example}
For some place prepositions like `behind, under\ldots', there are different forms based on whether they are used as standalones or as part of a noun unit:
\begin{example}
    \begin{tabular}{c c|c c}
        \textit{Noun unit} & \textit{Alone} & \textit{Lone usage} & \textit{Unit usage}\\\hline
        hinter & hinten & Ich stehe hinten & Ich stehe hinter dem Haus\\\hline
        unter & unten & Die Katze ist unten& Die Katze ist unter der Treppe
    \end{tabular}
\end{example}
Standalones are used when a general location  is specified: in other words when case $7$ is used without an associated noun. For example: \bt{ich stehe hinten} just means that I am standing somewhere behind, but a relative placement with respect to other noun landmarks is not specified. In this respect, one may also think of standalone prepositions as indefinite pronoun contractions: \bt{hinter etwas/hinten} 

We now consider \bt{da}-composites, which are a shorthand way to combine a preposition and \bt{das} ($\vm{P}[c,In]$) as a prefix in the form  \bt{da}. To form them, grab \bt{da} and a preposition of choice, and check if they chain a vowel (have consecutive vowels) in between:
\begin{itemize}
    \item \textit{Yes} - add an \bt{r} in between and smash them together
    \item \textit{No} - smash them together
\end{itemize}
\begin{example}
    \begin{tabular}{c c|c c|c c}
        \bt{dahin*} & towards there &\bt{darauf} & on that & \bt{dazu} & to that \\\hline \bt{daher*} & from there &\bt{damit} & with that& \bt{daraus} & from that
    \end{tabular}
\end{example}
Similarly, one may also combine \bt{was} in the form  \bt{wo} as a prefix with prepositions:
\begin{example}
    \begin{tabular}{c c|c c|c c}
        \bt{wohin*}& where to& \bt{worauf} & on what & \bt{wozu} & for what \\\hline \bt{woher*}& where from& \bt{womit} & with what & \bt{woraus}& from what
    \end{tabular}
\end{example}
You might notice that \bt{wozu} doesn't translate as expected. The truth is, although the given meanings are valid for such words, they are often not their only meanings. \bt{Damit} can also mean `so that', and \bt{dazu} can mean `for that purpose'. These double meanings have no pattern and must be learned over the course of time as is. Other combination variants not following the above rule like \bt{somit, infolge, zugunsten} also exist, sometimes with their own double meanings.

The entries with the asterisk (*) represent direct composition, which is a unique feature related to \bt{hin/her}. The \bt{da} or \bt{wo} in these entries are literal, and do not represent a contraction of \bt{das} or \bt{was}. Hence, their meanings deviate from  other examples, despite looking similar.

According to their meanings, different prepositions require  their associated nouns to be in certain cases. Often,  prepositions accept multiple cases, but some accept only one or two, which in rare cases may lead to forced casing contrary to expectation. First, a list of common prepositions:
\begin{setdef}
    \begin{tabular}{c|m{4cm}|c|m{5cm}}
        \textit{Verhältniswort}& \textit{Meaning}& \textit{Cases}& \textit{Sentence}\\\hline
        auf & \bt{on} & $2, 7$ & Ich klettere auf den Berg\\\hline
        von & \bt{from (location); by (passive)} & $5; 3$ &Der Zug fährt von Gleis drei ab\\\hline
        aus & \bt{from (boundary)} & $5$ & Sie kommt aus Deutschland\\\hline
        vor & \bt{before, in front of} & $2, 7$ & Er steht vor mir\\\hline
        in & \bt{in} & $2, 7$ & Wir gehen in die  Höhle\\\hline
        an & \bt{on, at} & $2, 7$ & Am(an + dem) Bahnhof\\\hline
        bei & \bt{at, in case of} &  $7$ & Sie blieb bei seiner Seite\\\hline
        nach & \bt{towards} & $4$ & Wir fahren nach Garmisch\\\hline
        zu & \bt{to, on, at} & $4, 7$ & Wir gehen zu Fuß\\\hline
        mit & \bt{with} & $3, 4$ & Ich komme mit\\\hline
        seit & \bt{since} & $5$ & Seit diesem Zeitalter\ldots\\\hline
    \end{tabular}\end{setdef}\begin{setdef}
    \begin{tabular}{c|m{4cm}|c|m{5.5cm}}
        für* & \bt{for} & $4[2]$ & Für das Reich!\\\hline
        %zugunsten & \bt{in favour of} & $6$ & Deren zugunsten\ldots\\\hline
        um* & \bt{around, at(time)} & $7[2]$& Um den Hals\ldots\\\hline
        mittels & \bt{by means of}& $6$& Mittels des Fahrrads\ldots\\\hline
        bezüglich & \bt{in consideration of (regarding)}& $6$& Bezüglich des Makels\ldots\\\hline
        jenseits & \bt{beyond the boundary of} & $6$ & Jenseits der Stadtgrenze\ldots
    \end{tabular}
\end{setdef}
Anytime the corresponding meaning needs to be described with an `of' at the end, case $6$ must be used. Hence, \bt{mittels} needs case $6$ despite technically being closer to case $3$. The ones marked with an asterisk (*) are exceptions, where the meaning and the [forced case] do not match. 

\bt{Aus} implies an exit from an enclosed boundary, whereas \bt{von} implies an exit from a certain starting point (like a person). \bt{Von} is also used as a case $1$ marker for a passive sentence, and a case $6$ marker for noun units without any additional features. For the second use, the case $6$ noun stays in \bt{Dativ}, since \bt{von} does not support \bt{Genitiv}. \bt{Aus} can also be used when describing the material composition of any substance: `made (out of) glass'

It is possible for \bt{zu} to be used as an emphasizer rather than a preposition: \bt{zu viel Energie} (too much energy). Here, \bt{Energie} is not in case $4$. Such uses are most likely impossible for any  preposition other than \bt{zu}, since I have never come across any other examples.

Finally, it is sometimes possible to place prepositions at the end of a noun unit. This is more common in specific fixed expressions and not generally done:
\begin{example}
    meiner Meinung nach, den Bach herunter\ldots
\end{example}
\section[Sentence]{Sentence (Satz)} \label{sec:destce}
Sentences generally follow a Subject-Verb-Object type construction similar to English, but with a catch.  In \autoref{sec:deverb}, it was mentioned that the CV always stays in second place, and the DC is kept at the end. This rule, which will be called the second place rule (SPR) is defined as:
\begin{syntax}
    \begin{itemize}
    \item For any verb chain in any sentence, the DC always stays at the very end, and the CV always stays exactly after the first unit in the sentence
    \item If any rule requires  the CV to be shifted to the end, the CV is appended to the end of the DC
    \end{itemize}
\end{syntax}
Note that the additional rules for $V\in\mathcal{S}$  apply when appending the CV to the end of the DC. Since this  is not an aspect change, the CV will never interact with anything in the DC, even for \bt{Futur 2}.

In the context of object verbs, it must be mentioned that \bt{sein, werden} when used as main verbs do not take case $2$ nouns, but rather two case $1$ nouns:
\begin{example}
    \begin{tabular}{c|c}
        \bt{Er($1$) ist ein kriegerischer Mann($1$)} & He($1$) is a belligerent man($2$)\\\hline
        \bt{Er($1$) wird sie($1$)}& He($1$) becomes her($2$)
    \end{tabular}
\end{example}
No other verbs I know do this. As to why this is the case  is because both of them make a \bt{Perfekt} using \bt{sein}, but require objects regardless. Since \bt{sein} as a helping verb only couples with non-object verbs, the necessary object must be put in case $1$ instead to preserve the rule.
\subsection[Punctuation]{Punctuation (Zeichensetzung)}
\textgerman{\bt{Meine Kämpfe gegen die Stadt trieben mich jenseits des Landes, und in den Bereich der Verrückten. Da fand ich einen Mann, der sagte \enquote{Heil der Sonne! Wir wären nichts ohne sie!} Verwirrt? So war ich auch. Auch zu sehen auf einem Schild waren die Wörter \enquote*{Mein Leid wird bald verbrannt werden.} Plötzlich klingelte eine Glocke, und Männer, Frauen, Ältere, Jungen; selbst die Tiere des Orts liefen langsam in Richtung eines Hauses. Ich wusste, dass irgendwas ernsthaft schiefgelaufen war : Alles fühlte sich an, als wäre es tot; wie die Vernunft schon geopfert worden war.}}

This paragraph has all the punctuation used. There are no usage differences compared to English, and you will recognise the minor symbolic changes easily. Did I need to write this nonsensical paragraph to showcase this? Not at all. But being unpredictable is a mark of a good villain. The above paragraph also contains many sentence based elements which have not been introduced yet, and gives a brief outlook of how  \bt{Deutsch} normally looks.

\subsection[Conjunction]{Conjunction (Verbindungswort)}
There are two types of relational conjunctions based on how they interact with the SPR:
\begin{enumerate}
    \item \textit{Shifting}: CV appended to the end of  DC.  Can always start sentences
    \item \textit{Normal}: CV at second position. Mostly cannot start sentences
\end{enumerate}
Note that unlike verb aspect variation, shifting conjunctions do not alter the form of the CV: they only throw it to the end. For uses with cell $\vm{T}(1,2)$, the prepositional argument (if any) simply gets slammed onto the CV: \bt{abfahren/abfährt}

If any relational conjunction starts a sentence, the consequent connecting clause has its CV immediately after the comma. This occurs because the entire previous clause along with its conjunction is treated as a single  placeholder  unit, so the SPR places the next CV immediately after it:
\begin{example}
        (Dass ich hier zu diesem Zeitpunkt lebe), ist ein Wunder
\end{example}
Additional conjunctions behave exactly as they would in English, and do not effect any special changes. 
\subsubsection{Dass}
That. Can be coupled with \bt{so} to give \bt{sodass} meaning `so that'. Shifting type conjunction, and is often ignored while speaking in order to avoid exiling the CV and to improve understanding, especially when speech fillers are used:
\begin{example}
    \begin{tabular}{m{7cm}|m{7cm}}
        \bt{Ich glaube, (dass) ich mag Honig} & I believe, (that) I like honey\\\hline
        \bt{Ich unterdrücke meine Wut, sodass ich in der Gesellschaft leben kann} & I suppress my fury, so that I in \ut{the} society \ut{live can}
    \end{tabular}
\end{example}
\subsubsection{Da/Weil/Denn}
Because. Among these, \bt{da/weil} are shifting, whereas \bt{denn} is normal. \bt{Denn} is unable to start sentences, whereas the other two have no such weakness. \bt{Da} is more formal than \bt{weil}. Also, this \bt{da} obviously does not mean `there' in this context.
\begin{example}
    \begin{tabular}{m{7cm}|m{7cm}}
        \bt{Wir müssen uns beeilen, weil der Zug gleich abfährt} & We must \ut{ourselves} hurry, because the train shortly \ut{departs}\\\hline
        \bt{Ich weiß  nicht alles, denn ich bin kein Gott} & I \ut{know not} everything, because I am \ut{no} god
    \end{tabular}
\end{example}
\subsubsection{Wenn, Falls}
If/when, In case. Both of them are shifting types, and mean slightly different things. \bt{Falls} is placed at the beginning of a  conditional argument, with the other argument expressing the precaution that has been taken against that situation, which would arise, were the condition to come true:
\begin{example}
    \begin{tabular}{m{6.5cm}|m{7cm}}
        \bt{Falls das Kind aufwacht, haben wir eine Überwachungskamera im Kinderzimmer} & In case the child wakes up, \ut{have we} a security camera in the children's bedroom\\\hline
        \bt{Wir haben einen Regenschirm dabei, falls es regnet} & We have an umbrella \ut{at this point}, in case it rains
    \end{tabular}
\end{example}
\bt{Wenn} on the other hand implies cause-consequence. Although it does mean `when' sometimes, such uses are rather context dependent and not that common, since `when' is normally written using \bt{als}. In the consequent solution, \bt{wenn} exclusively means `if'. 

\bt{Wenn} has an additional construction when it starts a sentence:
\begin{itemize}
    \item Eliminate \bt{wenn} and place the CV in its place
\end{itemize}
No other conjunction allows this.
\begin{example}
    \begin{tabular}{m{6.5cm}|m{7cm}}
        \bt{Spricht man über das Wetter, vergeudet er seine Zeit}& (If) \ut{speaks} one about the weather, is wasting \ut{he} his time\\\hline
        \bt{Wasser friert, wenn seine Temperatur unter Null sinkt} & Water freezes, if its temperature under zero \ut{sinks}
    \end{tabular}
\end{example}
Adverbs can also be used as relational conjunctions, which may either be normal or shifting.  The only difference between adverbial and pure conjunctions is that adverbials always count as a unit in the SPR, whereas pure conjunctions never do. Since there are too many adverbials, only pure conjunctions will be comprehensively  listed. First, \autoref{tab:destceT1} contains all elements  $E\in\mathcal{A}_\boxdot$, where anything that counts in the SPR is marked with the exponent $\mathbb{V}$.
\begin{table}[bth]
    \centering
    \begin{subtable}{\textwidth}
        \centering
    \begin{tabular}{|c|c|}
        \hline \textit{Conjunction} & \textit{Meaning}\\\hline
        \bt{und} & and\\\hline
        \bt{oder} & or\\\hline
        \bt{beziehungsweise} & respectively / more precisely\\\hline
        \bt{aber} & but\\\hline \bt{allein} & alone\\\hline
        \bt{sondern} & but rather\\\hline
        \bt{doch}& nevertheless / however\\\hline
        \bt{als} & as (only for comparisons)\\\hline
    \end{tabular}
    \caption{Singles: $\mathcal{A}_\boxdot[\mathbb{S}]$}
\end{subtable}
\begin{subtable}{\textwidth}
    \centering
    \begin{tabular}{|c|c|m{0.8cm}|}
        \hline \textit{Conjunction} & \textit{Meaning}& \textit{Max args.}\\\hline
        \bt{$\text{entweder}^{\mathbb{V}}$\ldots oder\ldots}& either\ldots or\ldots & $\infty$\\\hline
        \bt{$\text{weder}^{\mathbb{V}}$\ldots $\text{noch}^{\mathbb{V}}$\ldots}& neither\ldots nor\ldots & $\infty$\\\hline
        \bt{nicht nur\ldots sondern\ldots auch\ldots}& not only\ldots but also\ldots (surprising) & $2$\\\hline
        \bt{sowohl\ldots als auch\ldots}& not only\ldots but also\ldots (standard) & $2$\\\hline
        \bt{$\text{zwar}^{\mathbb{V}}$\ldots aber\ldots} & indeed\ldots but\ldots (statement + context) & $2$\\\hline
        \bt{$\text{einerseits}^{\mathbb{V}}$\ldots $\text{andererseits}^{\mathbb{V}}$\ldots} & on  one hand\ldots on the other hand\ldots & $2$\\\hline
        \bt{$\text{mal}^{\mathbb{V}}$\ldots $\text{mal}^{\mathbb{V}}$\ldots}& sometimes\ldots sometimes\ldots & $\infty$\\\hline
        \bt{$\text{teils}^{\mathbb{V}}$\ldots $\text{teils}^{\mathbb{V}}$\ldots}& partly\ldots partly\ldots & $\infty$\\\hline
    \end{tabular}
    \caption{Coupled: $\mathcal{A}_\boxdot[\mathbb{O}]$}
\end{subtable}
    \caption{Additionals: $\mathcal{A}_\boxdot$}
    \label{tab:destceT1}
\end{table}

The \bt{auch} in the three part conjunction \bt{nicht nur\ldots sondern\ldots auch\ldots.} stands before the second joined unit, but \bt{sondern} always stays at the very beginning of the second clause. Hence, if there are other remaining units, they collectively form the second argument, which is otherwise the null set:
\begin{example}
    Sie bringt (nicht nur) [Wein] (sondern) gibt den Kindern (auch) [Messer]
\end{example}
Hence, there are still only two arguments, and not three. \bt{mal\ldots, teils\ldots.} can be chained. 

At least all pure elements $E\in \mathcal{R}_\boxdot$  are in \autoref{tab:destceT2}. $\mathcal{R}_\boxdot$ also contains exactly two coupled conjunctions. Any normal adverbials  are marked with the exponent $\mathbb{N}$.
\begin{table}[bth]
    \centering
    \begin{tabular}{|c|c|}
        \hline\textit{Conjunction}& \textit{Meaning}\\\hline
        \bt{als} & when\\\hline \bt{ob} & if/whether (Yes/No Questions)\\\hline \bt{als ob}& as if (mostly \bt{K2})\\\hline
        \bt{damit} & so that / in order to \\\hline \bt{\{so\}dass}& \{so\} that\\\hline \bt{wenn} & if\\\hline \bt{weil/da} & because\\\hline
        \bt{falls} & in case \\\hline \bt{ohne dass} & without\\\hline \bt{dadurch, dass} & due to the fact that\\\hline \bt{indem}& using\\\hline
        \bt{obwohl / obschon / obzwar / obgleich} & although \\\hline \bt{$\text{denn}^{\mathbb{N}}$} & because\\\hline
        \bt{$\text{deshalb}^{\mathbb{N}}$ / $\text{deswegen}^{\mathbb{N}}$ / $\text{daher}^{\mathbb{N}}$ / $\text{darum}^{\mathbb{N}}$} & hence\\\hline
        \bt{$\text{trotzdem}^{\mathbb{N}}$ / $\text{dennoch}^{\mathbb{N}}$}& nevertheless\\\hline
        \bt{wohingegen} & against which / whereas \\\hline \bt{bevor / ehe} & before\\\hline \bt{bis}& until\\\hline \bt{während} & while\\\hline
        \bt{solange} & as long as\\\hline \bt{sobald}& as soon as\\\hline \bt{sooft} (pronounced \bt{so/oft})& as often as / whenever \\\hline
        \bt{nachdem} & after\\\hline \bt{seit / seitdem} & since\\\hline \bt{Je\ldots desto/umso\ldots} & the\ldots the\ldots (comparators)\\\hline
    \end{tabular}
    \caption{Relationals: $\mathcal{R}_\boxdot$}
    \label{tab:destceT2}
\end{table}

The comparators \bt{Je\ldots desto/umso\ldots.}  only take comparative adjectives as arguments:
\begin{example}
    \begin{tabular}{m{7cm}|m{7cm}}
        \bt{Je [länger] ich auf dieser Erde lebe, desto [stärker] werde ich} & The longer I live on this earth, the stronger I will become
    \end{tabular}
\end{example}
\bt{Dadurch, dass} literally means `through the fact that', and the first part \bt{dadurch} doesn't have any consistent placement rules, at least not to my knowledge. It need only exist somewhere in the first argument. To check the nature (shifting/normal) of any unknown adverbial, simply use \textit{dwds.de}: a resource linked in \autoref{sec:dermrk}.

Next  are  conjunction contractions related to \bt{dass} and \bt{damit}. These are only allowed when the case $1$ nouns of both argument clauses  are the same.
\subsubsection{Dass - zu + Grundform}
Works for anything  that uses \bt{dass} as a component. Eliminate the case $1$ noun and conjunction \bt{dass} from the affected clause, and use $\vm{T}(:,4)$, with the CV being pushed to the end:
\begin{syntax}
    DC + \bt{zu} + CV (\bt{Grundform})
\end{syntax}
The special rules pertaining to $\mathcal{S}$ still apply when the CV is exiled. \bt{Zu} always appears before the very last verb in the chain. This contraction  represents an aspect, which is applied to the time  carried over from the main clause:
\begin{example}
    \begin{tabular}{m{6.5cm}|m{6.5cm}}
        \textit{Expanded}& \textit{Contracted}\\\hline
        Der Mann ist  klug genug, dass er seinen Weg   herausfindet  & Der Mann ist klug genug, seinen Weg  herauszufinden \\\hline
        Er konnte nicht glauben, dass er verloren hat & Er konnte nicht glauben, verloren zu haben\\\hline
        Er konnte nicht glauben, dass er so viel  verlieren kann & Er konnte nicht glauben, so viel  verlieren zu können\\\hline
        Der Mann ist glücklich, dass er seinen Weg  hat herausfinden können & Der Mann ist glücklich, seinen Weg   herausgefunden haben zu können
    \end{tabular}
\end{example}
The argument with `that' either occurs along/after  its other argument, or is already complete before it. Hence, ONLY cells  $\vm{T}(:,2)$ can be used in the argument clause with \bt{dass}, which is why only two contraction types using $\vm{T}(:,4)$ exist and suffice. All cells can otherwise be used in the other argument. 

\textit{An exception} occurs with the $\vm{T}(2,4)$ type contraction if any $V\in\mathcal{S}$ is involved. In this case, the main verb is no longer permeable and interacts with the exiled $R\in\mathcal{X}$, causing a $\vm{T}(2,2)$ chain to be formed. Note that $R$ must follow the main verb due to this interaction according to $\vm{T}(2,2)$. All other rules apply as normal.

Sometimes, if the main clause has a small DC,  the contracted argument is placed before it as part of the sentence:
\begin{example}
    \begin{tabular}{m{7cm}|m{7cm}}
        \bt{Ich werde Besseres [mit meinem Geld anzufangen] wissen} & I will know to start something better with my money
    \end{tabular}
\end{example}
\subsubsection{Damit - um\ldots zu + Grundform}
Replaces the adverbial \bt{damit} meaning `in order to/so that'. The rules are exactly the same as the above \bt{zu + Grundform}, except that now an extra \bt{um} stands at the beginning of the contracted clause:
\begin{example}
    \begin{tabular}{m{7.5cm}|m{7.5cm}}
        \textit{Expanded}& \textit{Contracted}\\\hline
        Er verwendete die Segen der Dämonen, damit er sich aus der Schwierigkeit aufheben konnte& Er verwendete die Segen der Dämonen, um sich aus der Schwierigkeit aufheben zu können\\\hline
        Er zählt die Schafe, damit er die Winterzeit aushält & Er zählt die Schafe, um die Winterzeit auszuhalten
    \end{tabular}
\end{example}
$\vm{T}(2,4)$ based constructions are  rare for this contraction type due to its meaning.
\subsubsection{Ohne zu}
This shouldn't technically be here but there is no other fitting place for it:
\begin{setdef}
    \begin{tabular}{m{7.5cm}|m{6.5cm}}
        \textit{Rule 1} & \textit{Rule 2}\\\hline
        Verbs of perception (\bt{sehen, hören\ldots}) & \bt{bleiben, sitzen, stehen, liegen}\\\hline
        \bt{helfen} & Verbs of motion (\bt{gehen, fahren, kommen\ldots})\\\hline
        - & \bt{lernen, lehren, üben} 
    \end{tabular}
\end{setdef}
Under such usage, the aforementioned verbs take verb arguments according to the rules: 
\begin{enumerate} 
    \item Temporarily treated as members of $\mathcal{S}$ 
    \item Treated similar to $\mathcal{S}$, with the only difference being the \bt{Perfekt} aspect.  \bt{Perfekt}ion  occurs normally, and there are no double infinitive rules as otherwise observed with  members of $\mathcal{S}$
\end{enumerate}
Verb arguments serve as meaning supplements. There is unfortunately no good explanation as to why  $\mathcal{S}$ allows the second category of verbs only partial citizenship.
\begin{example}
    \begin{tabular}{m{6.5cm}|m{7cm}}
        \bt{Wir wollen hier stehen bleiben} & We want \ut{here standing stay}\\\hline
        \bt{Sie ist heute früher spazieren gegangen} & She has \ut{today earlier} \ut{walking gone}\\\hline
        \bt{Er hat ihn rennen sehen} & He has \ut{him} \ut{running seen}
    \end{tabular}
\end{example}
Original members from $\mathcal{S}$ can also be used on such chains, as illustrated in the first example.

\subsection[Interrogation]{Interrogation (Fragesatz)}
$\biguplus$ We start with the \textit{Yes/No} (\bt{Ja/Nein}) questions:
\begin{syntax}
    Switch the  case $1$  and CV units
\end{syntax}
Practically, this often means switching the first two units, so that the CV unit appears  first. Since conjunction usage is invalid in questions, the first clause of the SPR always holds, thereby simplifying this formation:
\begin{example}
    \begin{tabular}{c|c}
        \bt{Bist du tot?} & Are you dead?\\\hline
        \bt{Wollen wir uns wieder treffen?} & \ut{Want we us again meet?}\\\hline
        \bt{Ist er zu früh angekommen?} & Has he too early \ut{arrived}?
    \end{tabular}
\end{example}
The interrogative words (IWs) used for \textit{elaborations} are as follows:
\begin{setdef}
    \begin{tabular}{c|c|c|c}
    \bt{Wann} & When & \bt{Wie} & How\\\hline \bt{Wo} & Where & \bt{Was}& What\\\hline
    \bt{Warum}& Why &\bt{Wieso} & Why\\\hline \bt{Wer*} & Who & \bt{Welcher*}& Which
    \end{tabular}
\end{setdef}
The words with the * change form, while all others are invariant. The form variation of \bt{welcher} is exactly as in \autoref{tab:denounT9}, whereas \bt{wer} follows \bt{der} ($\vm{Q}(:,1)$). 

\bt{Warum} focuses on the underlying reason behind an occurrence, whereas \bt{wieso} asks how  the occurrence came to pass:
\begin{example}
    \begin{tabular}{c|c|c}
        \textit{Question}& \bt{Warum} & \bt{Wieso}\\\hline
        Why did you fall?& The floor was wet & I slipped on the wet floor
    \end{tabular}
\end{example}
Using one of these IWs as the first unit in a sentence declares an elaboration:
\begin{example}
    \begin{tabular}{c|c}
        \bt{Warum tust du das?}& Why are  doing \ut{you} this?\\\hline
        \bt{Woher kommst du?} & From where \bt{come} you?\\\hline
        \bt{Wohin willst du?} & \ut{To where} \ut{want} you (go)?
    \end{tabular}
\end{example}
The situation is slightly different with \bt{wie}, which often deals with scales and properties, thus requiring adverbs. Interestingly, \bt{wie} and the adverb  are treated as a single unit for the SPR:
\begin{example}
    \begin{tabular}{m{6.5cm}|m{7cm}}
    \bt{Wie spät ist es?} & How \ut{late} is it?\\\hline
    \bt{Wie bist du hierhergekommen?} & How have you come here?\\\hline
    \bt{Wie schnell  fahren wir?} & How fast are \ut{driving} we?
    \end{tabular}
\end{example}
As a reminder, adverbs do not have  gradations, so \bt{Wie schneller} is not allowed.
\subsection[Describer subclause]{Describer subclause (Beschreibungssatz)}
$\biguplus$ \autoref{tab:destceT3} represents $\vm{Q}$ with all  hook word forms. The syntax of German D-subclauses is exactly the same as English: Commas separate the subclause, which is introduced using a hook word, from the main sentence. If a preposition is used, it is placed exactly before the hook word. The second SPR clause is always applicable:
\begin{example}
    \begin{tabular}{m{7cm}|m{7cm}}
        \bt{Der Mann, der ein König geworden war, wird\ldots}& The man, who had  become king, is becoming\ldots\\\hline
        \bt{Das Kind, dessen Haustier allmählich sterben wird, war\ldots} & The child, whose pet will be dying gradually, was\ldots\\\hline
        \bt{Die Frau, deren Kind einen Wutanfall gemacht hatte, ist angekommen} & The woman, whose child had thrown a tantrum, has arrived
    \end{tabular}
\end{example}
$\vm{Q}$ can also be used as an emphasizing replacement for  $\vm{P}[a,In/De,3_p]$. This feature is formally introduced in \autoref{sec:despco},  and requires clear prior context for use:
\begin{example}
    \begin{tabular}{m{7cm}|m{7cm}}
        \bt{Einführung in elektrische Systeme und deren Anwendungen\ldots} & Introduction to electrical systems and their (the electrical systems specifically) applications\\\hline
        \bt{Einführung in elektrische Systeme und ihre Anwendungen} & Introduction to electrical systems and their (could be something mentioned earlier, maybe mechanical systems) applications
    \end{tabular}
\end{example}
Pronoun based subclauses require iteration of the pronoun, since $\vm{Q}$ strictly assumes $3_p$:
\begin{example}
    \begin{tabular}{m{7cm}|m{7cm}}
        \bt{Ich, der ich dieses Weltall erstellt habe, bin gelangweilt}& I, who \ut{this universe} \ut{created have}, am bored\\\hline
        \bt{Du, die du mit dem Kind gekämpft hast, bist trottelig} & You, who have \ut{with the child} \ut{fought}, are moronic
    \end{tabular}
\end{example}
The gender of $\vm{Q}$ is taken from the real gender of the person, similar to adjective based nominalizations.
\begin{table}[bth]
    \centering
    \begin{tabular}{|c|c|c|c|c|}
        \hline$\diamond$ & \bt{Männlich} & \bt{Weiblich} & \bt{Sächlich}& \bt{Mehrzahl}\\\hline
        \bt{Nom} & der & die & das & die\\\hline
        \bt{Akk} & den & die & das & die\\\hline
        \bt{Dat} & dem & der & dem & denen \\\hline
        \bt{Gen} & dessen & deren & dessen & deren\\\hline
    \end{tabular}
    \caption{Hook words: $\vm{Q}$}
    \label{tab:destceT3}
\end{table}

For indefinite descriptions, the IWs: \bt{was, welcher} are used with their corresponding form variations. \bt{Welcher} is used only for specific references in technical German, similar to the replacement possessives as described above. \bt{Was} is used as usual:
\begin{example}
    \begin{tabular}{m{6.5cm}|m{7cm}}
        \bt{Das Haus, welches sie gekauft haben, ist alt} & The house, which you \ut{bought have}, is old \\\hline
        \bt{Nichts, was wir machen, ist umsonst}& Nothing \ut{what} we do is \ut{for nothing}\\\hline
        \bt{Das Beste, was ich je gesehen habe} & The best \ut{what} I ever seen \ut{have}
    \end{tabular}
\end{example}
Contracted adjectives and sentence units also require \bt{was}, if further described. For the CV, \bt{was} is always treated as a third person singular, if used in case $1$. Under certain circumstances, \bt{warum, wo} can also be used as hook words, but these add  sentence units  rather than actual describer subclauses:
\begin{example}
    \begin{tabular}{m{7cm}|m{7cm}}
        \bt{Das Haus, wo ich wohne\ldots} & The house, where I live\ldots\\\hline
        \bt{Das Haus, in dem ich wohne\ldots}& The house, in which I live\\\hline
        \bt{Der Grund, warum ich nicht kommen kann\ldots} & The reason why I cannot come\ldots\\\hline
        \bt{Der Grund, wegen dessen ich nicht kommen kann\ldots} & The reason, because of which I cannot come
    \end{tabular}
\end{example}
\subsubsection{MdG Verkürzung}
The \bt{Mittelwort der Gegenwart (MdG)}, alternatively known as \bt{Partizip1} is an analogue to the English P1, which is formed by adding a \bt{-d} to the \bt{Grundform}. The sole exception is \bt{sein/seiend}, which doesn't end up mattering because it has no use case. None.
\begin{example}
    \begin{tabular}{m{7cm}|m{6cm}}
    Er fährt mit dem Bus, der hell und heiß brennt, zur Arbeit & Er fährt mit dem (hell und heiß brennenden) Bus zur Arbeit\\\hline
    Sie aß sein Bein, das schmerzhaft ausblutete, voller Freude & Sie aß sein (schmerzhaft ausblutendes) Bein voller Freude
    \end{tabular}
\end{example}
\bt{Voller Freude} is an old German relic representing case $6$, which means `full of joy'. Covered later in \autoref{sec:despco}. 

\bt{MdG} contractions can also be carried out for the \bt{sein + zu + Grundform} passive replacement. The CV is \bt{sein}, and \bt{zu + Grundform} immediately precedes it. For the contraction, \bt{sein} is eliminated, while the \bt{Grundform} becomes the \bt{MdG}:
\begin{example}
    \begin{tabular}{m{7cm}|m{7cm}}
        \textgerman{Wir werden morgen die frühen Prüfungen, die unter \enquote*{Altklausuren} zu finden sind, behandeln} & \textgerman{Wir werden morgen die frühen (unter \enquote*{Altklausuren} zu findenden) Prüfungen behandeln}
    \end{tabular}
\end{example}
There is no \bt{zu + MdV} counterpart to this contraction type.
\subsubsection{MdV Verkürzung}
The \bt{Mittelwort der Vergangenheit (MdV)}, alternatively known as \bt{Partizip2} is an analogue to the English P2, whose formation has been dealt with in \autoref{sec:deverb}.

This is the only point where the diligent reader may get confused, because traditional German grammar and I treat an \bt{MdV} expansion differently. Although I consider it as the \bt{Zustandsleidform} (state passive), grammar textbooks often use the \bt{Vorgangsleidform} (process passive) instead. Clearly, a justification is in order.

An embedded subclause acts identical to an adjective, in the sense that it describes a property of the associated noun. Much more important than the process of how this property came to be, is the fact that it exists. English agrees on this respect, and the consequent expansion is  logical and easier to grasp. The \bt{Vorgangsleidform} is quite counterintuitive and does not fit in many cases, since the `Ongoing' aspect references that the action is currently in progress at the time of reference. On the other hand, the aspects do not mean exactly what one may expect for the \bt{Zustandsleidform}, which will be clarified in the section on passives, and make the subclause interpretation intuitive.

Although one may argue that the \bt{Perfekt} aspect with process passives could be used, standard grammar does not allow this, and the rules surrounding subclause contractions are strict to begin with. Hence, any expansion of \bt{MdV} subclauses  will explicitly use the \bt{Zustandsleidform} in this solution, which does not follow standard texts:
\begin{example}
    \begin{tabular}{m{7cm}|m{7cm}}
        Die Frau, die trotz fehlender Beteiligung in der Masche getötet ist, ist schön & Die (trotz fehlender Beteiligung in der Masche getötete) Frau ist schön\\\hline
        Den Aufbau der Leistungselektronik, die in den verschiedenen Fahrzeugsteuergeräten eingesetzt waren, \ldots & Den Aufbau der (in den verschiedenen Fahrzeugsteuergeräten eingesetzten) Leistungselektronik\ldots
    \end{tabular}
\end{example}
Although the given examples do not reflect this, embedding is independent of other noun unit features: \bt{schnell fahrendes Auto / geschicktes Bild} without any articles is valid.

\ut{\bt{Warning}}: \bt{MdV} contractions cannot be  performed for any verb $V\in\mathcal{S}$.
\subsubsection{Miscellaneous}
\begin{enumerate}
\item The modal verbs \bt{Müssen, Können} can  be replaced using \bt{zu + Grundform}, which falls under the category  `Modal replacements'. Among them, \bt{müssen} may only be replaced in the active voice:
\begin{setdef}
    \begin{tabular}{m{4cm} c}
        \textit{Modal Verb} & \textit{Replacement}\\\hline
        Müssen  & haben + zu + Grundform \\\hline
        Können & Haupttätigkeitswort + [zu + Grundform][\bt{D-Subclause}]
    \end{tabular}
\end{setdef}
The first verb is the CV. Since  \bt{können}  uses $\vm{T}(1,4)$ for contraction, only $\vm{T}(1,:)$ is allowed in the subclause:
\begin{example}
    \begin{tabular}{m{7cm}|m{7cm}}
        \textit{Contracted} & \textit{Expanded}\\\hline
        Ich habe Brot zu schneiden & Ich muss Brot schneiden \\\hline
        Sie fand nichts zu sagen & Sie fand nichts, was sie sagen konnte\\\hline
        Ich habe ihn zu töten gehabt & Ich habe ihn töten müssen \\\hline
        Es gibt nur wenig zu machen & Es gibt nur wenig, was gemacht werden kann
    \end{tabular}
\end{example}
\item Both \bt{MdG, MdV} describer subclause contractions can be written in a different form, if a noun unit embedding is not desired. Simply eliminate the hook word and CV (assuming the DC is not a null set), and let the remaining phrase stay as is:
\begin{example}
    \begin{tabular}{m{7cm}|m{6.8cm}}
        \textit{Expanded} & \textit{Contracted}\\\hline
        Die Braut, die vor Glück strahlte, sah in die Kamera & Die Braut, vor Glück strahlend, sah in die Kamera\\\hline
        Ihr Vater, der von Wehmut ergriffen war, wandte sich ab & Ihr Vater, von Wehmut ergriffen, wandte sich ab
    \end{tabular}
\end{example}
This is especially useful if the noun is the name of a person. Sometimes, if the main sentence is small, the shortened subsentence can also be thrown to the very end, though this might cause confusion and should be avoided:
\begin{example}
    Ihr Vater wandte sich ab, von Wehmut ergriffen
\end{example}
\item Pronoun noun duplication is possible in general  for $\vm{P}[a,In,:](1,\mathbb{M},0\leftrightarrow\mathbb{H})$:
\begin{example}
    \begin{tabular}{c c|c c}
        \bt{Wir Reisenden\ldots} & We the travelling\ldots & \bt{Ihr Suchenden\ldots} & You the seeking\ldots\\\hline
        \bt{Wir Gefangenen\ldots} & We the prisoners\ldots & \bt{Sie Verletzten\ldots} & They the hurt\ldots\\\hline
        \bt{Sie Ungeheuer\ldots} & You monsters\ldots & \bt{Ihr Fahrer\ldots} & You drivers\ldots
    \end{tabular}
\end{example}
This may also  be done for some $\vm{P}[a,In,:](1,1,0\leftrightarrow\mathbb{H})$ under specific conditions: \bt{du Süße}
\item \bt{Wenn}-arguments using cell $\vm{T}(1,2)$ with \bt{man, es} as the case $1, 2$ nouns respectively may also be contracted. To do this, \bt{wenn, es} and \bt{man} are eliminated, and the CV is converted into its \bt{MdV} form. Sometimes the unit ordering can also be flipped:
\begin{example}
    \begin{tabular}{m{7cm}|m{7cm}}
        \textit{Contracted} & \textit{Expanded}\\\hline
        genau genommen & Wenn man es genau nimmt\ldots\\\hline
        verglichen mit [Zitronen] & Wenn man es mit [Zitronen] vergleicht\ldots\\\hline
        abgesehen davon & Wenn man es davon absieht\ldots\\\hline
    \end{tabular}\\
    anders ausgedrückt/gesagt, angenommen, gesetzt [in] den Fall, kurzgesagt, so gesehen/betrachtet, vorausgesetzt\ldots
\end{example}
\item Noun duplication  follows the typical \{\textit{CoM} content \textit{CoM}\} pattern:
\begin{example}
    Die Quelle [des Rheins], [des längsten Flusses innerhalb Deutschlands], liegt in der Schweiz
\end{example}
\end{enumerate}

\subsection[Order/Request]{Order/Request (Befehl/Aufforderung)}
\begin{setdef}
    \begin{tabular}{c|m{7cm}}
        Order cell & $\vm{T}(1,2)$ \\\hline
        Subsubsection naming & $\vm{S}[\vm{T}(1,2)](2_p,m,h)$, where ($m,h$) are variables\\\hline
        Grade keyword & \bt{Bitte}\\\hline
        Hypothesis based cmd & \bt{Möglichkeitsform (K2)} of \bt{können / haben}
    \end{tabular}
\end{setdef}
Compounding grade keywords and hypotheticals cannot cross into the next integral range, unless the grade is physically increased. \bt{Haben} hypotheticals are treated separately at the end.
\subsubsection{Cmd:(1,0)}
\begin{enumerate}
    \item The CV is kept as the first unit
    \item Eliminate any \bt{Umlaut} that appears due to irregular conjugation, and the \bt{-st} end. Sometimes, an additional \bt{-e} may appear. (Do not eliminate natural \bt{Umlaut})
    \item Eliminate \bt{du}, and add an \textit{ExC} (!), \textit{InQ} (?) or \textit{EoS} (.) at the end
    \item Add any grade keywords if desired
\end{enumerate}
Hypothesis commands follow interrogative structures, so only the last step  is applicable for them.
\begin{example}
    \begin{tabular}{m{6cm}|m{6cm}|c}
        \textit{Command} & \textit{Original/Translation} &\textit{Grade}\\\hline
        Könntest du bitte gehen? & \bt{Could you please go?} & $\approx0.8$\\\hline
        Töte die Ungerechtigkeit der Welt, bitte& Du tötest die Ungerechtigkeit der Welt & $0.5$\\\hline
        Zieh dich an! & Du ziehst dich an  & $-0.5$\\\hline
        Iss nicht normal & Du isst nicht normal & $0$
    \end{tabular}
\end{example}
Exceptions: \bt{sein/sei}, \bt{haben/hab}. 
\subsubsection{Cmd:(M,0)}
\begin{enumerate}
    \item The CV is kept as the first unit
    \item Eliminate \bt{ihr}, and add an \textit{ExC} (!), \textit{InQ} (?) or \textit{EoS} (.) at the end
    \item Add any grade keywords if desired
\end{enumerate}
\begin{example}
    \begin{tabular}{m{6cm}|m{6cm}|c}
        \textit{Command} & \textit{Original/Translation}& \textit{Grade}\\\hline
        Geht heim! & \bt{Go home!}& $-0.5$\\\hline
        Entschuldigt euch für dieses Verhalten & \bt{Excuse yourselves for this behaviour} & $0$\\\hline
        Könntet ihr leise sein? & \bt{Could you  quiet \ut{be}} &$0.5$
    \end{tabular}
\end{example}
\subsubsection{Cmd:(\{0,M\},H)}
\begin{enumerate}
    \item The CV is kept as the first unit and \bt{Sie} as the second
    \item Add an  \textit{InQ} (?) or \textit{EoS} (.) at the end
    \item Add any grade keywords if desired
\end{enumerate}
The highest politeness level does not accept exclamations, for obvious reasons.
\begin{example}
    \begin{tabular}{m{6cm}|m{6cm}|c}
        \textit{Command} & \textit{Original/Translation}& \textit{Grade}\\\hline
        Kommen Sie bitte herein & \bt{\ut{Come} \ut{You} \ut{please} inside} & $\mathbb{H}+0.5$\\\hline
        Entspannen Sie sich & \bt{Relax \ut{You} \ut{yourself}}& $\mathbb{H}$\\\hline
        Könnten Sie bitte einen Platz nehmen? & \bt{Could you please a seat \ut{take}?}& $\approx\mathbb{H}+0.8$
    \end{tabular}
\end{example}
Exception: \bt{sein/seien}

Rewriting \bt{Sie} based commands that do not employ hypotheticals can create signboard instructions. Simply use the second SPR clause and eliminate \bt{Sie}:
\begin{example}
    \begin{tabular}{m{7cm}|m{7cm}}
        \textit{Command} & \textit{Signboard}\\\hline
        Benutzen Sie (den) Aufzug im Brandfall nicht & Aufzug im Brandfall nicht benutzen\\\hline
        Schütteln Sie (den Saft) gut vor dem Öffnen & Vor dem Öffnen gut schütteln
    \end{tabular}
\end{example}
Often, the case $2$ unit is omitted or used without any additional features.
\subsubsection{Miscellaneous}
The \bt{haben K2} form is essentially a normal active voice \bt{K2} sentence in $1_p$. Without going into the technicalities of \bt{K2}, here are some examples. Note that this is still $\vm{T}(1,2)$:
\begin{example}
    \begin{tabular}{m{7cm}|m{7cm}}
        \bt{Ich hätte gern eine Banana} & I would like to have a banana\\\hline
        \bt{Wir hätten gern eine Fahrkarte} & We would like to have a ticket
    \end{tabular}
\end{example}
Mostly useful during shopping. The last remaining command form is the $1_p$ group commands, which follow the English structure using `let':
\begin{example}
    \begin{tabular}{c|c}
        \bt{Lass uns gehen!} & Let us go!\\\hline
        \bt{Los geht es!} & \ut{Off goes it}!
    \end{tabular}
\end{example}
The \bt{Cmd:(1,0)} form is applied on \bt{Lassen} for such group orders. \bt{Los geht's} is also commonly used as a German specific phrase.

\subsection[Hypothesis]{Hypothesis (K2/Möglichkeitsform)}
$\biguplus$ First, $\vm{T}_K$ must be defined. Only the aspect `Ongoing' exists, with the following distribution:
\begin{equation*}
    \vm{T}_K(1,1)=\vm{T}(2,2), \quad \vm{T}_K(1,2)=\vm{T}(1,2), \quad \vm{T}_K(1,3)=\vm{T}(1,3)
\end{equation*}
In order to apply the mode \bt{K2} on any CV and get its \bt{K2} root, the following steps are necessary:
\begin{enumerate}
    \item Take the \bt{Präteritum} root: $\vm{T}(1,1)$ of the CV
    \item If an \bt{Umlaut} can be added to an existent \bt{a/o/u}, add it
\end{enumerate}
Hence, $\vm{T}_K(1,1)\neq\vm{T}(1,1)$  to avoid overlap with $\vm{T}_K(1,2)$. That being said, the overlap is not completely absent, since a lot of verbs do not have any \bt{Umlaut}able letters, which makes the \bt{K2} and \bt{Präteritum} forms  coincide. Also, some verbs do not allow an \bt{Umlaut} addition despite having \bt{Umlaut}able letters.

To avoid this confusion, the $\vm{T}_K(1,3)$ form is often used as a substitute for $\vm{T}_K(1,2)$, as future assumptions are rare anyway. CV form variation follows the pattern in $\vm{T}(1,1)$, but with the \bt{K2} root.

The conjunction \bt{wenn} is used often when speaking of conditions, and is the only one which allows \bt{K2} usage within its arguments. The only exceptions appear for the sentence mode \bt{mittelbare Wirkart}, where \bt{dass, ob} among others may also take \bt{K2} arguments.

\bt{Wenn} is not strictly necessary for \bt{K2} usage - double usage  simply increases the uncertainty of the considered statements:
\begin{example}
    \begin{tabular}{m{7cm}|m{6.8cm}}
        \bt{Ich würde mich freuen, wenn er gehen würde}& I would \ut{myself} \ut{delight}, if he \ut{go} \ut{would}\\\hline
        \bt{Wenn ich [bloß/doch] ein Millionär wäre!} & If \ut{I} [only] a millionare \ut{were}!\\\hline
        \bt{Wären sie früher angekommen, wären wir früher gegangen} & Had \ut{they} earlier \ut{arrived}, would \ut{we} earlier \ut{have gone}\\\hline
        \bt{Wenn Sie  uns früher Bescheid gegeben hätten, hätten wir früher kommen können} & If you had earlier \ut{notice} \ut{given}, would \ut{we} \ut{earlier} have been able to come\\\hline
        \bt{Nähmest du einen Schritt weiter, sähest du seinen Tod} & If you take  a step further, you will see his death\\\hline
        \bt{Ginge die Sonne auf, käme die Vögel} & If the sun \ut{goes up}, the birds would come
    \end{tabular}
\end{example}
There are two special constructions associated with \bt{K2}:
\subsubsection{zu + als dass}
Too (something) to (something). Don't take  \bt{als dass}  word by word, it just means `to' collectively:
\begin{example}
    \begin{tabular}{m{7cm}|m{7cm}}
        \bt{Dieses Blut ist zu lecker, als dass ich es verpassen würde}  & This blood is too delicious to miss \\\hline
        \bt{Die Sanierung war zu teuer, als dass sie sich gerechnet hätte} & The renovation was too expensive to be cost effective
    \end{tabular}
\end{example}
The \bt{zu} in the first clause is necessary: it represents an exaggeration that due to the first condition, the second is highly improbable. The \bt{K2} time is matched according to the first clause.
\subsubsection{als ob / als wenn / als}
This is a direct analogue to `He behaves, as if\ldots' Sometimes, you can find \bt{K1} forms used for this purpose in written language:
\begin{example}
    \begin{tabular}{m{7cm}|m{6.5cm}}
        \bt{[Er tut so], als ob er Gott wäre} & [He pretends] as if he were god\\\hline
        \bt{[Du siehst aus], als hättest du tagelang durchgearbeitet} & [You look] like you have worked throughout the entire day
    \end{tabular}
\end{example}
Here, the \bt{K2} time is independent of the introductory clause.

\subsection[Voice]{Voice (Wirkart)}
$\biguplus$ The active voice will be termed \bt{Tätigform} while the passive as the \bt{Leidform}. The subject suppliers are:
\begin{setdef}
$\circledS = \begin{bmatrix}
    \texttt{von} & \texttt{durch}
\end{bmatrix}$
\end{setdef}
The process and state passives work exactly as described in \autoref{sec:enstce}. The process passive places emphasis on the fact that the action occurred, whereas the state passive is concerned only with the consequences of the action.
\subsubsection{Vorgangsleidform}
Also called the \bt{werden} ($V\in\mathcal{X}$) passive. To obtain its base cell:
\begin{enumerate}
    \item Send the \bt{MdV} form of the CV to the DC
    \item The new CV is \bt{werden}
    \item \bt{\ut{Exception}}: Iff  the \bt{Perfekt} form is necessary,  \bt{werden}  will become \bt{worden} in the DC instead of \bt{geworden}
\end{enumerate}
This creates $\vm{T}(1,2)$ for the process passive, which enables sentences like:
\begin{example}
    \begin{tabular}{m{7cm}|m{6cm}|c}
        \textit{Statement} & \textit{Translation}& \textit{Cell}\\\hline
        \bt{Das muss gemacht werden} & This must be done & $\vm{T}(1,2)$\\\hline
        \bt{Sein Feind ist  getötet worden} & His enemy has  been killed & $\vm{T}(2,2)$\\\hline
        \bt{Das Schwert wird gerostet worden sein müssen, als wir zurückkehren} & The sword will have to have been rusted, when  we return & $\vm{T}(2,3)$
    \end{tabular}
\end{example}
Note that since \bt{worden} is not a modal or main verb, it blocks any exiled $K\in\mathcal{X}$ from advancing leftwards beyond it, according to the second special rule for any $V\in\mathcal{S}$. 

By virtue of its meaning, the process passives for \bt{wollen} are formed by substituting it with \bt{sollen}:
\begin{example}
    \begin{tabular}{m{7cm}|m{7cm}}
        Man will einen neuen Kindergarten aufbauen & Ein neuer Kindergarten soll aufgebaut werden
    \end{tabular}
\end{example}
No other $V\in\mathcal{S}$ is substituted like this.

Although the process passive could technically be formed for reflexive usage, this is not grammatically allowed, and any occurrences are strictly colloquial. A frequently observed  use concerning the \bt{Vorgangsleidform} is the implication of $\vm{T}(1,3)$ despite using $\vm{T}(1,2)$:
\begin{example}
    \begin{tabular}{m{8cm}|m{6cm}}
        $\vm{T}(1,3)$ \textit{Statement} & \textit{Translation}\\\hline
        \bt{Widerrechtlich abgestellte Fahrzeuge werden abgeschleppt} & Illegally parked vehicles will be  towed\\\hline
        \bt{Die Gäste werden später begrüßt} & The guests will be greeted later
    \end{tabular}
\end{example}
Such examples depict situations where the action is fixed, rulebound or supposed to occur immediately in the near future. These conditions mimic the `simple' aspect  that English uses, which is why the best translation uses the English cell - $\vm{T}(1,3)$. However, since \bt{Deutsch} does not possess this aspect, the above method is a  workaround to create $\vm{T}(\texttt{simple},3)$. There exist no such workarounds for any other cells or sentence modes.

To reduce the DC length  when a $V\in\mathcal{S}$ is involved, there exist alternative constructions for  the \bt{Vorgangsleidform}.  However, not all $V\in\mathcal{S}$ can be replaced by any given method.

Since there are many such possible reformulations, only the 4 most important  have been listed below. This exclusion is inconsequential, since all of the omissions are based either on rule `1' or `2', and can replace either \bt{müssen} or \bt{können}:
\begin{setdef}
    \begin{tabular}{m{0.5cm}|m{6cm}|m{5cm}}
        $\clubsuit$ & \textit{Replacement}&  $V\in\mathcal{S}$\\\hline
        1 & Sein + zu + Grundform  & Müssen, Können, Sollen, Nicht dürfen\\\hline
        2 & Lassen sich + Grundform  & Können\\\hline
        3 & Sein + Eigenschaftswort(-bar/-lich) & Können\\\hline
        4 & Man  & Keine
    \end{tabular}
\end{setdef}
The fourth pattern is neither passive nor a replacement, but the performer is still unknown, and thus retains a passive-like meaning despite being in active voice. Examples for $\vm{T}(1,2)$:
\begin{example}
    \begin{tabular}{m{6.5cm}|m{7cm}}
        \bt{Replacement} & \bt{Original}\\\hline
        Der Bericht ist zu schreiben & Der Bericht [kann / muss / soll] geschrieben werden\\\hline
        Der Bericht ist nicht zu schreiben & Der Bericht darf nicht geschrieben werden\\\hline
        Die Gleichung lässt sich schöner schreiben & Die Gleichung kann schöner geschrieben werden\\\hline
        Die Schrift ist gut lesbar & Die Schrift kann gut gelesen werden\\\hline
        Man baut ein Haus & Ein Haus wird gebaut 
    \end{tabular}
\end{example}
The CV and $\vm{T}$-cell calculations in each case should be obvious.  \bt{Zu + Grundform} is  an inert component that remains at the end of the sentence (beginning of the DC) and doesn't participate in cell transformation.

Lastly, although it should be obvious, it must be mentioned that any passive replacement (except 4) cannot be subject to another passive voice change. So
\begin{example}
Der Bericht wird zu schreiben gewesen : $\vm{T}(1,2)$
\end{example}
is absolute nonsense. Since 4 is technically active voice, it can be passived, which just converts the sentence back into its original form.

\subsubsection{Zustandsleidform}
Also called the \bt{sein} ($V\in\mathcal{X}$) passive. To obtain its base cell:
\begin{enumerate}
    \item Send the \bt{MdV} form of the CV to the DC
    \item The new CV is \bt{sein}
\end{enumerate}
The aspects of the \bt{Zustandsleidform} have a slightly different interpretation than normal. Since the state passive focuses on the consequence of an action, $\vm{T}(1,:)$ means that the consequence of an action currently exists. $\vm{T}(2,:)$ implies that the consequence existed, but no longer exists. Since the state passive declares  an  absolute and guaranteed state, no $V\in\mathcal{S}$ can be used with it.
\begin{example}
    \begin{tabular}{m{7cm}|m{7cm}}
        \bt{Die Tür ist geschlossen} & The door is closed\\\hline
        \bt{Die Tür ist geschlossen gewesen} & The door has been being closed
    \end{tabular}
\end{example}
The second example means that the door being closed was a condition that existed, but has since stopped existing. Although both perspectives seem to deal with this similarly, English is focusing on the action, while \bt{Deutsch} is looking at its consequence. $\vm{T}(2,:)$ for the \bt{Zustandsleidform} is practically never used.
\subsection[Subjective usage]{Subjective usage (Wahrscheinlichkeitsgebrauch)}
$\biguplus$ A verb $V$ may be used subjectively  under specific conditions:
\begin{enumerate}
    \item $V\in$ defined probability markers
    \item Special $\vm{T}_S$-cell forms are employed
\end{enumerate}
Such a usage attaches the  probability associated with $V$ to the state/action described by the sentence. The possible markers along with their approximate probabilities are:
\begin{setdef}
    \begin{tabular}{c|c|c|c}
        \textit{V} & \textit{Probability (\%)} & \textit{V} & \textit{Probability (\%)}\\\hline
        Müssen & $99$&  Müssen (K2) & $90$\\\hline
        Dürfen (K2) & $70$&   Können & $60$\\\hline 
        Können (K2) & $\leq50$ & Mögen & $\leq50$
    \end{tabular}
\end{setdef}
According to the surity of the declaration, the appropriate marker should be used. There are only two real $\vm{T}_S$-cells:
    \begin{equation*}
\vm{T}_S(1,1)=\vm{T}(2,2), \quad \vm{T}_S(1,2)=\vm{T}(1,2), \quad \vm{T}_S^*(1,3)=\vm{T}_K(1,3) 
    \end{equation*}
No future forms exist in this pattern, so any future uncertainties are declared using $\vm{T}_K(1,3)$.  Applying the subjective sentence mode occurs as follows:
\begin{enumerate}
    \item Exile the CV to the DC
    \item The new CV is $V$, which always uses its $\vm{T}(1,2)/\vm{T}_K(1,2)$ forms (depending on if it requires \bt{K2})
\end{enumerate}
Note that probability markers $\notin\mathcal{S}$, and should be treated  separately. Hence, they can also be applied on the \bt{Zustandsleidform}, even though $\mathcal{S}$ usage is illegal with it:
\begin{example}
    \begin{tabular}{m{8cm}|m{5.5cm}}
        \bt{Er könnte haben kommen müssen} & He might have had to come\\\hline
        \bt{Das mag in diesem Fall stimmen} & That may be true in this case\\\hline
        \bt{Die Tür muss geschlossen sein} & The door must be closed\\\hline
        \bt{Die Tür muss aufgebrochen worden sein} & The door must have been broken open
    \end{tabular}
\end{example}

\subsection[Speech]{Speech (Redeweise)}
$\biguplus$ \bt{Deutsch} allows unit reordering, unlike English:
\begin{example}
    \begin{tabular}{m{7cm}|m{7cm}}
        \multicolumn{2}{c}{\textit{Direct speech}}\\\hline
        \textgerman{\bt{Der Lehrer sagte, \enquote{Ihr müsst die Aufgabe heute vervollständigen.}}} & The teacher said, \enquote{You must complete the task today.}\\\hline
        \textgerman{\bt{\enquote{Aus dem Weg!}, schrie er}} & He screamed, \enquote{Out of the way!}
    \end{tabular}
\end{example}
\subsubsection{Mittelbare Redeweise}
The IACs for indirect speech:
\begin{setdef}
    \begin{tabular}{c|c|c}
        \textit{Nature/Mode} & \textit{IAC} & \textit{Comments}\\\hline
        \bt{Aussage} & \bt{dass} & -\\\hline
        \bt{Fragesatz (Ja/Nein)} & \bt{ob} & Opener always first\\\hline
        \bt{Fragesatz (Erläuterung)} & - & The question remains a sentence unit\\\hline
        \bt{Befehl} & \bt{dass} & Often contracted 
    \end{tabular}
\end{setdef}
The explicit sentence mode `Speech' corresponds to \bt{Konjunktiv 1 (K1)}, for which  $\vm{T}\neq\vm{T}_R$.

The diligent reader would have realized by now that any CV  has only two types of form variations, corresponding to cells $\vm{T}(1,2)$ and $\vm{T}(1,1)$.  Hence, any sentence mode must have  defined CV forms for both these cells to have the `complete' range at its disposal. But for some reason, the sentence mode \bt{K1}  has no defined $\vm{T}_R(1,1)$ forms, which makes it impossible to create $\vm{T}_R(:,1)$.

To circumvent this, all of the statements made about the past relative to the opener are grouped into $\vm{T}_R(2,2)$ or the \bt{Perfekt} aspect. This does introduce ambiguity about   past statements, but fatal errors are rather rare. 

$\vm{T}_R(1,2)$ is based on $\vm{T}(1,2)$ with the conditions:
\begin{enumerate}
    \item The \bt{K1} root is the VR
    \item $(3_p,1)$ follows $(1_p,1)$
    \item An additional \bt{e} is \ut{always} added at the end of the \bt{K1} root in $(2_p,:,0)$
    \item If the \bt{Grundform} ends in \bt{-eln}, then the \bt{e} is eliminated in the \bt{K1} root
    \item \bt{Grundform} in \autoref{tab:deverbT2} implies `VR + \bt{en}' irrespective of whether the infinitive end was \bt{-n/-en}
\end{enumerate}
The \bt{K1} root knows no  irregularity. So \bt{sein} becomes \bt{sei/seien}, making for an incredibly simple form variation\ldots until you realize that half of these forms are exactly the same as  $\vm{T}(1,2)$. This is especially a problem for regular verbs, since it defeats the purpose of \bt{K1}, leaving the interpretation up to context. To compensate for this, \bt{K2} is used instead but with a \bt{K1} meaning, which is a strange oversight in an otherwise structured perspective. Since this \bt{K2} has nothing to do with hypotheticals, it will be renamed to \bt{K3} instead.

\bt{K3} forms are exactly the same as \bt{K2}, which means that only the three cells of $\vm{T}_K$ exist for use, further compressing the expressive range. Despite the drop in accuracy, one can still get by:
\begin{example}
    \begin{tabular}{m{7cm}|m{7cm}}
        \bt{Der Lehrer befahl uns,  unsere Aufgabe zu vervollständigen} & The teacher ordered us  to complete our task\\\hline
        \bt{Er forderte, dass wir aus dem Weg gingen} & He demanded, that we  get out of the way\\\hline
        \bt{Er erklärt, (dass) er sei müde} & He is declaring that he is tired\\\hline
        \bt{Er fragte uns, ob wir frei haben} & He asked us if we have free time\\\hline
        \bt{Er fragte uns, wann wir frei haben} & He asked us when we have free time
    \end{tabular}
\end{example}

Indirect arguments in the replacement  forms for \bt{Aussage} do NOT use the optional mode \bt{K1}. Prepositions  placed before the speaker unit in the opener will be marked $\mathfrak{A}$, while those placed after it $\mathfrak{E}$:
\begin{setdef}
    \begin{tabular}{c|c|c|c}
        \textit{Preposition} & \textit{Placement} & \textit{Meaning} & \textit{Surity}\\\hline
        \bt{laut} &  $\mathfrak{A}$ & according to & Very High\\\hline
        \bt{gemäß} & $\mathfrak{E}, \mathfrak{A}$ & in accordance with & High\\\hline
        \bt{nach} & $\mathfrak{A}$ & according to & Medium\\\hline
        \bt{zufolge} &  $\mathfrak{E}$ & in consequence to & Full
    \end{tabular}
\end{setdef}
The English counterparts have been provided solely to give an idea for why all of these require the speaker unit to use the \bt{Dativ} case. One may also indicate the perceived certainty of  the reported statements using the appropriate preposition. 

If \bt{wie}[as]($\mathfrak{A}$) is employed, the speaker unit must  stay in case $1$:
\begin{example}
\begin{tabular}{m{6cm}|m{7cm}}
\textit{Contracted} & \textit{Original}\\\hline
[Gemäß einer Forschung] gehen\ldots & Die Forschung sagt, dass\ldots gehen\\\hline
[Wie verschiedene Zeugen bestätigten], zeigte die Ampel rot & [Verschiedene Zeugen bestätigten], dass die Ampel rot gezeigt habe\\\hline
[Den Angaben zufolge], ist er tot & [Die Angaben ausdrücken], dass er tot sei
\end{tabular}
\end{example}
\subsubsection{Miscellaneous}
The verbs \bt{sollen, wollen} can be used to recount the words of another without confirming their accuracy:
\begin{example}
    \begin{tabular}{m{6.5cm}|m{7cm}}
    \bt{Er soll ein erfolgreicher Anwalt sein} & [According to rumours/I have heard that] he is a successful lawyer\\\hline
    \bt{Er will ein erfolgreicher Anwalt sein} & [He claims/Supposedly] he is a successful lawyer
    \end{tabular}
\end{example}
The [text in square brackets] gives the  meaning encoded by \bt{wollen/sollen} in the respective cases. The  problem:  there are absolutely no VC differences with any regular uses that mean `want/should'. Instead, contextual clues do the heavy lifting, which is once again a strange oversight for an otherwise structured perspective.

Using  \bt{sollen, wollen} in this manner makes it possible  to  mix them with the \bt{Zustandsleidform}.

$\biguplus$ \bt{K1} is also used for the `Blessing' mode. Due to limited usage and the reuse of \bt{K1} forms, it does not deserve a separate subsection in \bt{Deutsch}:
\begin{example}
    \begin{tabular}{m{5cm}|m{5cm}|m{4cm}}
    \bt{Möge Gott euch segnen} & - & May god \ut{you} bless \\\hline
    \bt{Es lebe der König!} & \bt{Möge der König leben!} & May the king live!\\\hline
    \bt{Er fahre zur Hölle} & \bt{Möge er zur Hölle fahren} & May he go to hell\\\hline
    \bt{Gott helfe uns!} & \bt{Möge Gott uns helfen!} & May god help us!
    \end{tabular}
\end{example}
All the above sentences can  be universally written using \bt{mögen} based orders, but sometimes they are rephrased to sound natural/impactful.  \bt{Mögen} uses its subjective meaning in such usage.

The `Blessing' mode may also be used for generating indirect speech for a request:
\begin{example}
    \begin{tabular}{m{7cm}|m{8cm}}
        \bt{Der Kellner sagt, (dass) sie möge bitte Platz nehmen} & The waiter is requesting  her to take a seat\\\hline
        \bt{Die Abgeordneten baten, (dass) \{man\} möge ihrem Antrag zustimmen} & The representatives \ut{were requesting} that \{people\} consent to their application
    \end{tabular}
\end{example}

In technical \bt{Deutsch},  multiple considerations are often declared about something without hypothesizing: for example as the starting point for a proof, or derivation of a result, where it is more about logical reasoning and argumentation rather than unconfirmed statements. Such statements along `assumptions, argumentation, or instructions without a commanding tone'  use  \bt{K1}:
\begin{example}
    \begin{tabular}{m{5.5cm}|m{8cm}}
        \bt{Gegeben sei das skizzierte Rahmentragwerk} & \ut{Given is the sketched truss}\\\hline
        \bt{Betrachtet werde das Energiebordnetz eines Brennstoffzellenfahrzeugs} & The (on-board electrical system) of a (fuel cell automobile) is \ut{being} considered\\\hline
        \multicolumn{2}{m{13.5cm}}{\bt{Der Polizist rief dem Radfahrer zu, (dass) er  sofort anhalten müsse. [Er dürfe nicht auf die Autobahn fahren]} - [Instruction]}
    \end{tabular}
\end{example}
Since an assumption is to be declared, and the actor is not important, such structures are restricted to $\vm{T}_R(1,2)$ with the passive voice (unless dealing with instructions), and require the DC to stay as the first unit of the sentence. Such a technical restructuring is unique to \bt{Deutsch}.

\subsection[Negativity]{Negativity (Verneinung)}
$\vm{P}[b,In]$ based:
\begin{example}
    niemand, nichts, keiner, wenig, nirgendwo\ldots
\end{example}
The matrix $\vm{N}$ is given by:
\begin{setdef}
        $\vm{N} = \begin{bmatrix} \texttt{kein} & \texttt{nicht} \end{bmatrix}$
\end{setdef}
Noun unit negations are performed using \bt{kein}, which has been introduced as the negative article in \autoref{sec:denoun}:
\begin{example}
    \begin{tabular}{m{7cm}|m{7cm}}
    \bt{Das ist kein Spiel} & This is no game\\\hline
    \bt{Hier ist kein Geld zu verdienen} & There is no money to be earned here
    \end{tabular}
\end{example}
VC negations use  \bt{nicht}, which in general stands as  the first word in the unit it negates. If the meaning of the entire sentence is to be inverted, it stays at the very end, but before the DC:
\begin{example}
    \begin{tabular}{m{7cm}|m{7cm}}
        \bt{Sie haben mir das nicht erzählt} & You \ut{have} not told me this\\\hline
        \bt{So kannst du nicht überleben} & You can\ut{ }not survive like this\\\hline
        \multicolumn{2}{c}{\bt{Mengenangabe}}\\\hline
        \bt{Nicht jede Fruchtsorte bringt die gleichen Vorteile} & Not every  \ut{fruit type} brings the same advantages\\\hline
        \bt{Nicht alle verhalten sich vernünftig} & Not \ut{all} behave rationally
    \end{tabular}
\end{example}
\subsection[Timekeeping]{Timekeeping (Zeitmessung)}
$\mathfrak{R}24$:
\begin{syntax}
    (Hour) \bt{Uhr} (Minutes)
\end{syntax}
\begin{example}
    \begin{tabular}{c|c}
        acht Uhr fünfundfünfzig & 08:55\\\hline siebzehn Uhr zweiundvierzig & 17:42
    \end{tabular}
\end{example}
Markers for $\mathfrak{R}12$:
\begin{setdef}
    \begin{tabular}{c|c|c|c}
        \bt{halb} & half to (not past) & \bt{viertel} & quarter\\\hline
        \bt{vor} & before & \bt{nach} & after\\\hline
        \multicolumn{4}{c}{Uncertainty markers}\\\hline
        \bt{kurz} & shortly ($\approx5$min) & \bt{gleich / fast} & around / almost\\\hline
        \multicolumn{4}{c}{Half markers (HM)}\\\hline
        \bt{morgens} & AM & \bt{abends} & PM
    \end{tabular}
\end{setdef}
Only \bt{kurz} can be combined with \bt{vor, nach}. The reference time is  the future by default, with  time preposition  \bt{um}:
\begin{example}
    \begin{tabular}{m{6cm}|m{6cm}|c}
        \bt{Es ist kurz [vor] halb vier} & It \ut{is} shortly half past three & $\approxeq3:30-\epsilon$\\\hline
        \bt{viertel nach eins morgens} & quarter past one in the morning & $1:15$ AM\\\hline
        \bt{Es ist viertel vor zwei} & It is quarter to two & $1:45$\\\hline
        \bt{(kurz nach) zehn vor acht} & (shortly after) ten minutes to eight & $7:50 (+\epsilon)$\\\hline
        \bt{Um fast drei viertel acht} & at around \ut{three quarters after eight} & $8:45$
    \end{tabular}
\end{example}
The standard time question is: 
\begin{setdef}
\bt{Wie spät ist es?} : \ut{How late is it?}
\end{setdef}
Calendar units have the following substitutes:
\begin{setdef}
    \begin{tabular}{c|c|c|c}
        \bt{vor} & before & \bt{nach} & after\\\hline
        \multicolumn{4}{c}{Uncertainty markers}\\\hline
        \bt{einige Tage} & Some days & \bt{ungefähr} & approximately
    \end{tabular}
\end{setdef}
The date preposition is \bt{am = an + dem}, since every single month and date is grammatically \bt{männlich}:
\begin{example}
    \begin{tabular}{m{7cm}|m{7cm}}
        \bt{Am zwanzigsten Februar 2019\ldots} & On the twentieth February 2019\ldots\\\hline
        \bt{Der dritte Oktober ist der Tag der deutschen Einheit} & Third October is the day of german unity
    \end{tabular}
\end{example}
The standard date question is:
\begin{setdef}
    \bt{Welches Datum / Welcher Tag ist heute?} : Which date / day is today?
\end{setdef}
\section[Specific constructions]{Specific constructions (Besondere Aufbauten)} \label{sec:despco}
\subsubsection{Blackletter script (Fraktur)}
The German blackletter script, or \bt{Fraktur} is an old $16^{th}$ century script that is often found in equally old texts,  on street signs in \bt{Deutschland}, or in traditional businesses. \bt{Fraktur} is no longer used and has been phased out, but has its own angular charm. If one desires, one may seek:
\begin{example}
    \begin{tabular}{m{7cm}|m{7cm}}
    $\mathfrak{Das}$ $\mathfrak{ist}$ $\mathfrak{Fraktur}.$ $\mathfrak{Aber}$ $\mathfrak{ob}$ $\mathfrak{die}$ $\mathfrak{Schrift}$ $\mathfrak{noch}$ $\mathfrak{es}$ $\mathfrak{wert}$ $\mathfrak{ist?}$ & Das ist Fraktur. Aber ob die Schrift noch es wert ist?
    \end{tabular}
\end{example}
\subsubsection{Case Mirroring (Fallspiegelung)}
When \bt{als, wie} are used as comparators, the case of their second argument mirrors that of their first:
\begin{example}
    \begin{tabular}{m{7cm}|m{7cm}}
        \bt{Er($1$) war so groß wie sie($1$)} & He was as big as her\\\hline
        \bt{Besser dem Kind($4$) etwas schenken als dem Mann($4$)} & Better to the child \ut{something gift} than to the man
    \end{tabular}
\end{example}
English does  this for all cases except $1$.
\subsubsection{Odd words (Seltsame Wörter)}
Some adjectives  look like \bt{MdV} or \bt{MdG} forms of  verbs that don't exist:
\begin{example}
    \begin{tabular}{c c c|c c c}
        \multicolumn{3}{c}{\textit{MdV}} & \multicolumn{3}{|c}{\textit{MdG}}\\\hline
        verzwickt & berühmt & berüchtigt & anwesend & weitgehend & abwesend
    \end{tabular}
\end{example}
The meanings of some \bt{MdV} and \bt{MdG} adjectives have  no obvious relation  to the verbs they are derived from:
\begin{example}
    \begin{tabular}{c c|c c}
        \multicolumn{2}{c|}{\textit{MdV}} & \multicolumn{2}{c}{\textit{MdG}}\\\hline
        gelassen & \bt{serene} & dringend & \bt{urgent}\\\hline
        besessen & \bt{obsessed} & umgehend & \bt{immediately}
    \end{tabular}
\end{example}
The superlative phrase `One of the most\ldots' can take any general adjective, or even stay without it:
\begin{example}
    \begin{tabular}{m{14cm}}
    Einer der großen/größeren/größesten Schriftsteller Deutschlands\ldots\\\hline
    Eines seiner Versehen ist dafür Schuld
    \end{tabular}
\end{example}
\subsection[Forced/Permitted actions]{Gezwungene/Gestattete Handlungen}
Marker verbs are employed for both types:
\begin{setdef}
        $\vm{D} = \begin{bmatrix} \texttt{lassen} & \texttt{lassen} \end{bmatrix}$
\end{setdef}
The distinction between forced and permitted actions occurs via context or reformulation, since both are marked by \bt{lassen}. The diligent reader will question immediately the fact that no $E\in\vm{D}$ can $\in\mathcal{S}$, which contradicts the definition made in \autoref{sec:deverb}. This objection, while correct, does not affect any of the explanations made until this point.

\bt{Lassen} has only three possible use cases:
\begin{enumerate}
    \item Main verb without a verb argument
    \item Passive replacement: \bt{sich lassen}
    \item Forced/Permitted actions, which need a verb argument
\end{enumerate}
The irregularities that deny any $E$ from having complete membership of $\mathcal{S}$ lie solely in the passive voice. As a main verb without an argument, everything works as normal, while passive replacements cannot be subject to another passive voice change. Hence, the only exceptional configurations exhibited by \bt{lassen} occur for forced/permitted actions. Add to this the fact that such sentences normally do not employ the passive voice,  makes $\texttt{lassen}\in\mathcal{S}$ hold for most practical applications.

The sole irregularity is that \bt{lassen} uses its \bt{MdV} form when passived, but not when \bt{Perfekt}ed if it takes a verb argument:
\begin{example}
    \begin{tabular}{m{6.5cm}|m{6.5cm}}
        \bt{Ich muss mir meine Weisheitszähne ziehen lassen} & I must \ut{from myself my wisdom teeth pull let}\\\hline
        \bt{Jemand hat seine Münze fallen lassen} & Someone has his coin \ut{fall let}\\\hline
        \bt{Ich werde lernen gelassen} & I am being made to learn\\\hline
        \bt{Er lässt mich lernen} & He is making me learn\\\hline
        \bt{Lass mich gehen!} & Let me go!
    \end{tabular}
\end{example}
\subsection[Cases Akkusativ, Dativ]{Fälle Akkusativ, Dativ}
There exists a unique interaction between these cases when motion of any sort is involved, which has been hinted at in \autoref{sec:deprep} while explaining the forms: \bt{hinter/hinten, unter/unten\ldots}. The destination of any such movement  always has an associated preposition, and stays in \bt{Akk}($2$), assuming that the  preposition supports this. If not, the required case of the preposition takes over.

If a location of rest is declared without motion, \bt{Dat}($7$) is used for it. Any time phrases always use \bt{Dat}($7$) irrespective of motion:
\begin{example}
    \begin{tabular}{m{7cm}|m{7cm}}
        \textit{Akkusativ} & \textit{Dativ}\\\hline
        Sie hängt den Mantel an die Wand & Der Mantel hängt an der Wand\\\hline
        Sie fährt in die Stadt & Sie fährt nach der Stadt\\\hline
        \multicolumn{2}{c}{Zeit Dativ}\\\hline
        Regen beginnt in zwei Stunden & Ich bin vor vier Jahren angekommen
    \end{tabular}
\end{example}
Although this rule has no exceptions,  it may sometimes seem like $7$  is used despite possible interpreted motion:
\begin{example}
    Wir treffen uns am(an + dem) Bahnhof
\end{example}
This is not an exception, since the movement target is assumed as the other person, not the \bt{Bahnhof} itself.
\subsection[Contractions]{Verkürzungen}
These contractions are different from the ones considered earlier, since there is no reformulation  involved, and are merely methods of compressing  words  for ease of speech:
\begin{itemize}
\item Some articles can be pushed into prepositions
\item If the contracted word is \bt{es}, the \textit{EoL}(') is  employed
\item The \bt{ei} at the beginning of any indefinite article can be replaced by the \textit{EoL}
\end{itemize}
\begin{setdef}
    \begin{tabular}{c c|c c}
        für/das & fürs & zu/dem & zum*\\\hline
        zu/der & zur & an/dem & am\\\hline
        an/das & ans & auf/das & aufs\\\hline
        von/dem & vom & hinter/dem & hinterm\\\hline\hline
        \multicolumn{2}{c}{Los geht es!} & \multicolumn{2}{|c}{Los geht's!}\\\hline\hline
        einen & 'nen & einem & 'nem
    \end{tabular}
\end{setdef}
\bt{Zum*} and \bt{zudem} mean entirely different things, which makes this a case of forced contraction. Any other contraction is otherwise not absolutely necessary.
Not every article and preposition can be pushed together, which is why  the above table has been made comprehensive with respect to this rule.

\ut{Fun fact}: Some prepositions like \bt{außer}(apart from)($5$) were used so many times with neutral/male noun singulars, that \bt{außerdem}(in addition to) became a conjunction over time.
\subsubsection{Word elimination (Wortbeseitigung)}
In two specific cases, it is possible to eliminate a word and form a  unit from what remains:
\begin{enumerate}
    \item \bt{Sein} (as the CV for $\vm{T}(1,2)$ in the \bt{Zustandsleidform}) : \bt{Betreten (ist) verboten}
    \item \bt{Von : Eine Menge (von) Leute}
\end{enumerate}
The unit formed by method 1 cannot be used in sentences, and must stand alone. In contrast, method 2 forms a typical noun unit.
\subsubsection{Allgemeine Nennwortverkürzungen mittels Hakenwörter}
Hook words ($\vm{Q}$) can  replace nouns if their previous context is known, and emphasize their connection with any consequent usage. Such an emphasis attempts  to clarify the link between noun units in a paragraph, especially useful if multiple similar nouns are dealt with at once. An example for this type of replacement has been provided in the subsection `Describer subclause' with technical German.
\begin{example}
    \begin{tabular}{m{14cm}}
        Der Mann mag klein sein, aber [dessen] Wirkung ist groß\\\hline
        Mein Vater ist verschwunden. Ich kann [den] nicht mehr sehen\\\hline
        Der Gefangene ist zäh. [Dem] entkommt keine Auskunft
    \end{tabular}
\end{example}
Every such use for $\vm{Q}(1\leftrightarrow3,:)$ has an equivalent  $\vm{P}[a,In,3_p](1\leftrightarrow3,:)$ that can replace it. For $\vm{Q}_i(4,j)$, the equivalence is with $\vm{P}_i[a,De,3_p\{j\}](:,:)$, for any one choice of $i\in\{M,F,N\}$ and $j\in\{1,\mathbb{M}\}$ at a time.
\subsection[Infinitive units]{Grundformeinheiten}
$\vm{T}(:,4)$ stays at the end, with no replacement forms necessary  for any $V\in\mathcal{S}$:
\begin{example}
    \begin{tabular}{m{7cm}|m{7cm}}
        \bt{[Die vermissten Kinder zu suchen] hat  Vorrang} & \bt{*}[To be searching for the missing children] \ut{has} priority\\\hline
        \bt{[Ohne Wasser Arbeit gesucht zu haben] ist  ungewöhnlich} & [To have searched for work without water] is strange\\\hline
        \bt{[In dieser Welt leben zu müssen] wirkt manchmal nervig} & [To have to live in this world] is sometimes irritating\\\hline
        \bt{[Nicht mehr gepflegt zu sein] kann gefährlich wirken} & [To not be taken care of] can be dangerous
    \end{tabular}
\end{example}
$\vm{T}(1,4)$ infinitive units can also be  used  without  \bt{zu}. Such formulations may resemble reformulated polite commands, but are completely separate from them. $\vm{T}(2,4)$ units without \bt{zu} stand in for dRAs of \bt{weil/nachdem}:
\begin{example}
    [Unordentlich geworden], sah er tot aus : Da/Nachdem er [] war, sah er tot aus\\ Ein Unfall ist beim [Nägel schneiden] passiert
\end{example}
All infinitive units are always \bt{sächlich} and \bt{einzahlig}.

\bt{*}Keep in mind that since \bt{Deutsch} does not possess the \texttt{simple} aspect, $\vm{T}_{de}(1,4)$ corresponds to  $\vm{T}_{en}(1\leftrightarrow2,4)$ in English, making both translations possible depending on context.
\subsection[Fillers]{Füllwörter}
\begin{itemize}
    \item \bt{Halt}: just. Not to be confused with the command form of \bt{halten}:\begin{example}
        Ich habe halt zu viel um meine Ohren \end{example}
    \item \bt{Eben}: just (in the context of time):\begin{example}
        Er ist eben angekommen \end{example}
    \item \bt{Ich glaube}: I think/believe
    \item \bt{Doch}:  As a conjunction `nevertheless'. As a filler, it is a stress additive without having a concrete meaning. Can be used to answer questions where the expected answer is clear from the tone and context, so \bt{doch} becomes a placeholder that answers the other person through the tone\begin{example}
        A : Wir müssen bald Geld verdienen. B : Doch (\bt{Of course})\\ Wenn ich doch mehr Geld hätte! - \bt{If [only] I had more money!}\end{example}
    \item \bt{Bitte}: Filler similar to \bt{doch}\begin{example}
        A : Meine Mutter, sie ist\ldots \\ B : Wie bitte? (\bt{What did you say?})\\ A : Ich habe gesagt, meine Mutter ist eine Maschine geworden \\ B : Bitte? (\bt{Excuse me?})\\ A : Ja, sie gehört nun den Maschinengeistern\\ B : Bitte? (\bt{What are you on?}) \end{example}
    \item \bt{Klar}: obviously/clearly/naturally
    \item \bt{Mal}: The noun meaning is `time' as in `many times, one time', but the filler meaning is varied. It mostly carries a meaning of `for the moment' or `this time' \begin{example} Probiere das mal aus \end{example}
    \item \bt{Na}: `Well, what do you think?' Behaves like \bt{doch}, but  elicits a response from the other person. It  can also be used as an informal greeting along  `Yo, what's up?' \begin{example} Na, sollen wir gehen? \end{example}
    \item \bt{Oder/ne}: Seeks a reaction to the statement at the end of which they are placed \begin{example}Du hast kein Leben, ne? \end{example}
\end{itemize}
\subsection[It]{Es}
\bt{Es} has some other uses outside of being $\vm{P}_N[a,In,3_p](1\leftrightarrow2,1)$:
\begin{enumerate}
\item Can refer to  sentence units, similar to what pronouns ($\vm{P}[a/b/c,In]$) do for noun units. References are  made to the action described in the sentence unit rather than to the entire unit:\begin{example}\begin{tabular}{m{9cm}|m{5cm}}
    Sie las ein Buch, und ich tat es auch & Ein Buch lesen\\\hline
    Die anderen waren müde. Sie war es nicht & müde sein \end{tabular}\end{example}
\item \bt{Deutsch} does not allow the case $2$ unit for an object verb to be a main sentence directly: one must  introduce such a case $2$ unit as the dRA of \bt{dass}. \bt{Es} must consequently be added as a placeholder in the ndRA as its expected object. \bt{Es} is not used if the object verb always requires main sentence case $2$ units: \begin{example}\begin{tabular}{m{14cm}}
        Ich bedauere es, dass ich nicht kommen kann\\\hline Er hat es übernommen, allen Bescheid zu geben \\\hline Ich hoffe (es), dass er reich wird \end{tabular}\end{example}
    This \bt{es} disappears if the order of dRA/ndRA is switched.
\item It is possible in general to switch the position of the units for cases $1,2$ in any main sentence (not including the dRA of a shifting relational) without consequence, assuming they are distinct entities. Since non-object verbs have no defined case $2$ unit, this switch either uses units from other cases for the purpose (if they exist), or generates an \bt{es}:\begin{example}\begin{tabular}{m{7cm}|m{7cm}}
    Ein Unfall ist an der Kreuzung geschehen & An der Kreuzung ist ein Unfall geschehen\\\hline
    Ein Unfall ist geschehen & Es ist ein Unfall geschehen \end{tabular}\end{example}
    Sometimes, this switch can use a generated \bt{es} despite the presence of other units, but the validity of such a switch is no longer guaranteed:\begin{example}
        \textcolor{red}{Es erschrickt sie vor jeder Maus} : Es ist ein Unfall an der Kreuzung geschehen \end{example}
\item If the case $1$ unit of any sentence unit is an unknown non-living thing/a natural phenomenon, \bt{es} is used to denote it:\begin{example}\begin{tabular}{c|c}
    Es klingelte plötzlich & \bt{It was ringing suddenly} \\\hline Es nieselt seit Stunden & \bt{It is drizzling since hours}\\\hline Es ist sehr spät & \bt{It is very late} \\\hline Es ist mir kalt & \bt{It is cold \ut{in me}}\end{tabular}\end{example}
    A similar use is possible when speaking about a general topic. Most of these are fixed expressions:\begin{example}\begin{tabular}{c|c}
        Es handelt sich um ($2$) & \bt{It \ut{operates itself around}} : \bt{It is about}\\\hline 
        Es hängt von ($5$) ab & \bt{It \ut{hangs from}} : \bt{It depends on}\end{tabular}\end{example}
        A preposition with its associated argument always exists in such uses. If the argument is a sentence unit, it is declared after the main sentence with a \textit{CoM} separation, and the preposition is replaced with its \bt{da}-composite (\autoref{sec:deprep}):\begin{example} Es hängt davon ab, [wie er plant, weiterzugehen]\\ Es handelt sich darum, [warum die Vögel verschwunden sind] \end{example}
\end{enumerate}
\subsubsection{Unpersönliche Leidform}
Sentences using object verbs [without an object / with non-case $2$ objects] employ [\bt{es} / the non-case $2$ object] as  the  case $1$ unit in their \bt{Leidform}:\begin{example}\begin{tabular}{c|c}
    Viele Menschen feiern & Es wird gefeiert\\\hline Sie dachten & Es war gedacht \\\hline Der Schulleiter dankte dem Förderverein & Dem Förderverein wurde gedankt\\\hline Das hilft niemandem & Nimandem ist (damit) geholfen
\end{tabular}\end{example}
In some rare circumstances, non-object verbs can form the \bt{Vorgangsleidform} in \bt{Deutsch}. I have no idea why this is allowed:
\begin{example}\begin{tabular}{c|c}
        Sie lachen & Es wird gelacht\\\hline Sie klatschen & Es wird geklatscht
\end{tabular}\end{example}
The \bt{Zustandsleidform} does not display this idiosyncrasy.

\subsection[Ancient relics]{Altdeutsch}
\begin{german}
Wenn du diese Sätze schwer verstehst, bedeutet das, du bist für diesen Schritt noch nicht bereit.
\begin{enumerate}
    \item Wesfall Eigenschaftswörter: Manchmal, anstatt den Wesfall ($6$) auf ein Nennwort anzuwenden, setzt man  ein laut $\vm{G}(4,2)$ gebeugtes Eigenschaftswort in deren Nennworteinheit ein. Solche Eigenschaftswörter  bleiben immer unverändert, unabhängig vom Stand  des verbundenen Nennworts \begin{example}
        mit voller Genuss, Erlanger Kirche, Bamberger Dom, Nürnberger Rathaus \end{example}
        Bezüglich dieser Regel sind auch einige alte Bauweisen seltens zu beobachten, Überbleibsel aus einer alten Zeit, die ihre alten Bedeutungen verloren haben, und drücken heute etwas ganz anders aus: \begin{example} \begin{tabular}{m{14cm}}Eines Tages : An diesem willkürlichen Tag, der lediglich als  Ausgangspunkt für diese Geschichte betrachtet wird\\\hline Eines Nachts : An einer beliebigen Nacht, die jetzt als Ausgangspunkt erachtet wird \end{tabular}\end{example}
    \item Redewendungen: Wenn ein Satzteil die bisher erläuterten grammatischen Regeln willkürlich durchbricht,  zusammen nicht passende Wörter besitzt, oder eine unerwartete Bedeutung aufweist,  ist er höchstwahrscheinlich  eine Redewendung \begin{example} \begin{tabular}{c}
        Kommt Zeit, kommt Rat : \bt{Things will automatically get better with time}\\\hline Im Eifer des Gefechts : \bt{In the heat of the moment} \end{tabular}\end{example}
    \item Es gibt: Eine Floskel aus einem alten Bau, in dem das Tätigkeitswort \enquote{geben} Dasein darstellte. Die Floskel ändert sich auf keinen Fall, und braucht dazu eine Fall $2$ Nennworteinheit
\end{enumerate}
\end{german}
\section[Closing remarks]{Closing remarks (Abspann)} \label{sec:dermrk}
\subsection[References]{References (Verweise)}
\subsubsection{Grammar}
\begin{setdef}
    \begin{tabular}{c|c}
        \hline\textit{\bt{Resource}} & \textit{\bt{Type}}\\\hline
        \href{https://www.youtube.com/@LearnGermanOriginal}{Learn German} & YouTube channel (YTC)\\\hline
        \href{https://www.youtube.com/@MrLAntrim}{Herr Antrim} & YTC\\\hline
        Chat GPT & AI\\\hline
        Deutsch intensiv - Grammatik C1 (Angelika Fülleman) & Book\\\hline
        \href{https://www.youtube.com/@yourgermanteacher}{Your German teacher} & YTC\\\hline
    \end{tabular}
\end{setdef}
\subsubsection{Useful websites}
\begin{setdef}
    \begin{tabular}{m{5cm}|m{9cm}}
        \hline\textit{\bt{Website}} & \textit{\bt{Description}}\\\hline
        \href{https://www.dwds.de}{Digitales Wörterbuch der deutschen Sprache} & Comprehensive online library of all that is German, but only available in German. This is the most authentic source for meanings, $\vm{C}$ and $\vm{T}$, but requires good structural understanding for use\\\hline
        \href{https://conjugator.reverso.net/conjugation-german.html}{Reverso verb conjugator} & Exercise caution with helping verbs for $\vm{T}(2,2)$, sometimes incorrect. Useful for picking up synonyms\\\hline
        \href{https://der-die-das-train.com/}{Der die das Train} & $\vm{C}$ along with explanations pertaining to \autoref{sec:denoun}\\\hline
        \href{https://www.verbformen.de/de-en/}{Netzverb Wörterbuch} & $\vm{C}, \vm{T}$ along with meanings. This is the best English counterpart to \textit{dwds} and tailor made for learners, but I never saw its charm\\\hline
    \end{tabular}
\end{setdef}
\subsubsection{Exposure}
\begin{setdef}
    \begin{tabular}{c|m{9cm}}
        \hline\textit{\bt{Resource}} & \textit{\bt{Description}}\\\hline
        ARD Bibliothek & Financed by the radio tax (\bt{Rundfunkgebühr}), and does not work outside Germany\\\hline
        \textcolor{red}{FAU Kurse} &  Die Methode der finiten Elemente, Elektrische Energiespeichersysteme, Elektrifizierung von Fahrzeugen und Flugzeugen. You likely cannot access these\\\hline
        \href{https://www.pushpak.de/}{Pushpak} & Die Geschichten der Götter, schlicht übersetzt. Es gibt auch Hörbücher\\\hline
        YouTube & The algorithm  answers those who dare to ask. Try \href{https://www.youtube.com/@Testkanal.v1}{testkanal} for starters\\\hline
        Deutschlandradio & Financed by the radio tax (\bt{Rundfunkgebühr}), and does not work outside Germany\\\hline
        \href{https://www.youtube.com/@vorlesungen-am1erlangen477}{Applied mathematics} & Engineering mathematics in German\\\hline
    \end{tabular}
\end{setdef}
\subsubsection{Supporters}
\begin{setdef}
    \begin{tabular}{c|m{9cm}}
        \hline\textit{\bt{Artist}} & \textit{\bt{Art}}\\\hline
        Black Throne & \href{https://www.youtube.com/watch?v=EcTWY93HQgI}{Shin'ainaru Teki}\\\hline
        White Bat Audio & Too many to list. This man is legendary\\\hline
        Neo Stahl & \href{https://www.youtube.com/watch?v=e43F2jDoqH8&t=1335s}{Neon Checkpoint}, \href{https://www.youtube.com/watch?v=C_YSqw6hPSg&t=3592s}{Heavy Ordinance}\\\hline
        Shamisen Mugen Beats & \href{https://www.youtube.com/watch?v=M_2tG0_PL74}{V1}, \href{https://www.youtube.com/watch?v=Xz6KkS47OR4}{V2}, \href{https://www.youtube.com/watch?v=PseiObSfS-s}{V3}\\\hline
        Example$^{\mathsf{Ex}}$ & Example\\\hline
        Ex2$^{\mathfrak{E}}$ & Ex2\\\hline
    \end{tabular}
\end{setdef}
\subsection[Authors]{Authors (Schriftsteller/[innen])}
All authors should choose a sign from the math library (except \verb|$\boldsymbol$|) and mark any additions they make to \autoref{sec:dermrk} with it.
\begin{syntax}
    \bt{Main author} : \textsanskrit{ओजस अमित सेवेकर}
\end{syntax}
\begin{setdef}
    \begin{tabular}{cc|cc}
        \hline\bt{\textit{Name}} & \bt{\textit{Sign}} & \bt{\textit{Name}} & \bt{\textit{Sign}}\\\hline
        Test example & $\mathsf{Ex}$ & Test ex2 & $\mathfrak{E}$ \\\hline
    \end{tabular}
\end{setdef}